% !TEX program = xelatex
% 电工电子技术 · 第一章第一讲（Chapter one-1）· 精讲讲义（lecture-deepdive 规范）
\documentclass[11pt]{ctexart}
%=====================================================================
% theme-violet.tex —— 讲义精讲 · 紫花主题
% 用法：\documentclass[11pt]{ctexart}
%       \input{assets/theme-violet}
%       \begin{document} ... \end{document}
% 只用 ctexart 自带能力 + xcolor/enumitem/fancyhdr/pifont/amsmath，
% 不依赖 tikz/tcolorbox，保证 MiKTeX 可直接编译。
%=====================================================================

\usepackage[a4paper,left=2.0cm,right=2.0cm,top=2.1cm,bottom=2.3cm,%
            headsep=12pt,footskip=22pt]{geometry}
\usepackage{amsmath,amssymb}
\usepackage{xcolor}
\usepackage{array}
\usepackage{enumitem}
\usepackage{pifont}
\usepackage{fancyhdr}
\usepackage[hidelinks]{hyperref}

%---------------------------------------------------------------- 配色
\definecolor{violet}{RGB}{107,78,143}      % 主色：紫
\definecolor{violetdeep}{RGB}{63,44,92}    % 深紫：标题/强调
\definecolor{petal}{RGB}{183,159,208}      % 花瓣：装饰线、花饰
\definecolor{bandviolet}{RGB}{247,243,251} % 浅紫底：档案/提示框
\definecolor{bandcream}{RGB}{252,249,242}  % 米白底：易错/注意框
\definecolor{rulepurple}{RGB}{217,207,230} % 分隔线
\definecolor{inkgray}{RGB}{92,96,110}      % 次要文字

%---------------------------------------------------------------- 字体
% 有 Noto CJK 用 Noto，没有就退回 ctex 默认（Windows 下为宋体/黑体）
\IfFontExistsTF{Noto Serif CJK SC}{\setCJKmainfont{Noto Serif CJK SC}}{}
\IfFontExistsTF{Noto Sans CJK SC}{\setCJKsansfont{Noto Sans CJK SC}}{}

%---------------------------------------------------------------- 版式
\setlist{nosep,leftmargin=1.9em}
\linespread{1.28}
\setlength{\parskip}{3.5pt}
\setlength{\parindent}{0pt}
\setlength{\tabcolsep}{5pt}
\renewcommand{\arraystretch}{1.28}

%---------------------------------------------------------------- 花饰
% \ding{95} = 五瓣小花，\ding{96} = 星形花；都来自 pifont，必定存在
\newcommand{\orn}{\textcolor{petal}{\ding{95}}}          % 行内小标
\newcommand{\ornstar}{\textcolor{petal}{\ding{96}}}      % 星形花
\newcommand{\flowerrule}{%                                % 居中花饰分隔线
  \par\vspace{7pt}%
  \noindent\makebox[\linewidth]{\color{petal}%
    \rule{0.20\linewidth}{0.6pt}\hspace{8pt}{\small\ding{95}}\hspace{4pt}%
    \raisebox{1.1pt}{\tiny\ding{95}}\hspace{4pt}{\small\ding{95}}\hspace{8pt}%
    \rule{0.20\linewidth}{0.6pt}}%
  \par\vspace{7pt}}

%---------------------------------------------------------------- 标题与节头
\newcommand{\doctitle}[3]{%   #1 主标题  #2 副标题  #3 元信息（可用 \\ 分行）
  \begin{center}
    {\color{petal}\small \orn\enspace #2\enspace \ornstar}\\[7pt]
    {\LARGE\bfseries\color{violetdeep} #1}\\[10pt]
    {\color{petal}\rule{0.55\linewidth}{1pt}}\\[4pt]
    {\small\color{inkgray} #3}
  \end{center}\vspace{2pt}}

\newcommand{\hdr}[2]{%   #1 章节号（如 §0）  #2 标题
  \par\vspace{12pt}%
  \noindent\colorbox{violetdeep}{\color{white}\bfseries\small\;#1\;}\hspace{7pt}%
  {\color{violetdeep}\large\bfseries #2}\par\nobreak
  \vspace{3pt}{\color{petal}\hrule height 1.2pt}\vspace{6pt}}

\newcommand{\subhdr}[1]{%
  \par\vspace{6pt}%
  \noindent{\color{violet}\bfseries #1}\par\vspace{1pt}}

%---------------------------------------------------------------- 提示框
\newcommand{\callout}[1]{% 浅紫底：讲义档案、读法说明
  \par\medskip\noindent
  \fcolorbox{rulepurple}{bandviolet}{\parbox{\dimexpr\linewidth-2\fboxsep-2\fboxrule}{\vspace{2pt}#1\vspace{2pt}}}%
  \par\medskip}

\newcommand{\warnbox}[1]{% 米白底：易错、边界条件
  \par\medskip\noindent
  \fcolorbox{rulepurple}{bandcream}{\parbox{\dimexpr\linewidth-2\fboxsep-2\fboxrule}{\vspace{2pt}#1\vspace{2pt}}}%
  \par\medskip}

\newcommand{\gapitem}[2]{% #1 缺口类型（紫） #2 内容
  \par\vspace{2pt}\noindent
  \hangindent=1.4em\hangafter=1
  {\color{petal}\small\ding{96}}\enspace{\bfseries\color{violetdeep}#1}\enspace #2\par}

% 前提旁注：只放在「需要理解的那一步」后面，给陈述 + 本题代入，不重证。
\newcommand{\aside}[1]{%
  \par\vspace{1pt}\noindent
  \hangindent=1.6em\hangafter=1
  {\small\color{inkgray}%
    \hspace*{0.6em}{\color{petal}$\triangleright$}\hspace{2pt}#1}%
  \par\vspace{2pt}}

%---------------------------------------------------------------- 表格
% 本机 MiKTeX 的 colortbl 与 array 版本不匹配（\rowcolor 会报错），
% 故表格统一用「上粗紫线 + 表头 + 中细线」的三线式，不依赖 colortbl。
\newcommand{\tabtop}{\par\vspace{3pt}{\color{violet}\hrule height 1.2pt}\vspace{1pt}}
\newcommand{\tabmid}{\par\vspace{1pt}{\color{petal}\hrule height 0.5pt}\vspace{1pt}}
\newcommand{\tabbot}{\par\vspace{3pt}{\color{rulepurple}\hrule height 0.5pt}\vspace{8pt}}
\newcommand{\thead}[1]{{\bfseries\color{violetdeep}#1}}

%---------------------------------------------------------------- 页眉页脚
\pagestyle{fancy}
\fancyhf{}
\renewcommand{\headrulewidth}{0.4pt}
\renewcommand{\headrule}{\hbox to\headwidth{\color{rulepurple}\leaders\hrule height \headrulewidth\hfill}}
\fancyhead[L]{\small\color{inkgray}\songti \doctag}
\fancyhead[R]{\small\color{inkgray}\songti \coursetag}
\fancyfoot[C]{\small\color{inkgray}\orn\enspace{\color{violet}\thepage}\enspace\orn}

% 每篇讲义在导言区设这两个值
\newcommand{\doctag}{}
\newcommand{\coursetag}{}
\newcommand{\settags}[2]{\renewcommand{\doctag}{#1}\renewcommand{\coursetag}{#2}}

%---------------------------------------------------------------- 数学记号
\newcommand{\vect}[1]{\boldsymbol{#1}}
\newcommand{\dd}{\,\mathrm{d}}
\newcommand{\ddx}{\frac{\mathrm{d}}{\mathrm{d}x}}
\addtolength{\jot}{2pt}



\settags{C1-1 · 精讲}{电工电子技术}
\usepackage{tikz}
\usetikzlibrary{arrows.meta,calc,decorations.pathmorphing}

% ------- 电路图统一样式（纯 tikz，本机无 circuitikz）-------
\tikzset{
  cwire/.style={thick},
  cres/.style={fill=white,draw=black,thick},          % 电阻：空心矩形
  csrc/.style={fill=white,draw=black,thick},           % 源：空心圆
  carr/.style={->,>=Stealth,thick},
  varrow/.style={->,>=Stealth,thick},
  cnode/.style={fill=black},
  clab/.style={font=\small},
  cval/.style={font=\small,inner sep=1pt}
}
\newcommand{\resbox}[3]{% #1 x  #2 y  #3 标签
  \draw[cres] (#1-0.28,#2-0.75) rectangle (#1+0.28,#2+0.75);
  \node[cval] at (#1+0.5,#2) {#3};}
\newcommand{\resboxH}[3]{
  \draw[cres] (#1-0.75,#2-0.28) rectangle (#1+0.75,#2+0.28);
  \node[cval] at (#1,#2+0.5) {#3};}
\newcommand{\vsrc}[3]{% #1 x #2 y #3 标签；+ 在上
  \draw[csrc] (#1,#2) circle (0.34);
  \node[cval] at (#1,#2+0.55) {$+$};
  \node[cval] at (#1,#2-0.55) {$-$};
  \node[cval] at (#1-0.55,#2) {#3};}
\newcommand{\isrc}[3]{% 电流源：圆内一横，箭头向上
  \draw[csrc] (#1,#2) circle (0.34);
  \draw[thick] (#1-0.34,#2) -- (#1+0.34,#2);
  \node[cval] at (#1-0.6,#2) {#3};}

\begin{document}

\doctitle{直流电路的第一层语言：参考点、参考方向与电源模型}
{C1-1 逐页精讲 · 从「电路是什么」到「电源为什么可以等效互换」}
{依据：\texttt{Chapter one -1-2026中白上传.pdf}（p1--p60，共 60 页；本讲义页码一律指该册 PDF 物理页）\ \textperiodcentered\ 2026-10-23}

%=============================================================
\hdr{§0}{课件坐标、记号与三处拼接}

\subhdr{0.1 本册实际讲了哪几节，各占哪几页}

先说一件必须纠正的事：有一种流传的说法把 p24--p60 整段都算作 1.4.2 参考方向，
逐页核对页码后这不成立。p26 起页眉已经换成 \textbf{1.5 Ideal circuit elements}，
p44 起换成 \textbf{1.6 Practical sources}，p53 起是例题 E1。
实际分布如下（本讲义页码一律指本册 PDF 物理页）：

\par\vspace{2pt}
\noindent\begin{tabular}{@{}p{1.3cm}p{1.4cm}p{13.0cm}@{}}
\tabtop
\thead{页码} & \thead{节号} & \thead{页眉原文：本页内容 $\longrightarrow$ 本讲义对应}
\\
\tabmid
p1 & --- & Contents：全课程七章目录
\\
p2--p4 & 1.1 & Function and components：本讲纲目、1.1 的四项内容 $\to$ §1
\\
p5 & 1.1.1 & Definition of the circuit：电路＝电流流过的闭合路径 $\to$ §1.1
\\
p6--p8 & 1.1.2 & Types and Functions：电力电路／信号电路 $\to$ §1.2
\\
p9--p10 & 1.1.3 & Basic circuit components：电源、负载、中间环节、负载增大 $\to$ §1.3
\\
p11 & 1.1 & 四个提问（全页无答案）$\to$ §1.3
\\
p12 & 1.2 & Basic quantities：六个基本量的复习清单 $\to$ §1.5
\\
p13--p16 & 1.3 & Circuit conditions：闭合、开路、短路三态 $\to$ §1.4
\\
p17--p20 & 1.4.1 & Reference point：参考点、电位、$V_A$ 例题两问 $\to$ §2
\\
p21--p25 & 1.4.2 & Reference direction：参考方向三写法、Example 1／2／3 $\to$ §3、§4
\\
p26--p32 & 1.5.1 & Ideal passive element：电阻、电感、电容、忆阻器 $\to$ §4.2
\\
p33--p36 & 1.5.2 & Ideal active element：理想电压源、理想电流源 $\to$ §6.1、§6.2
\\
p37--p43 & 1.5.3 & Power absorbed／supplied：功率判据与四张 3A／3V 图 $\to$ §5
\\
p44--p45 & 1.6.1／\newline 1.6.2 & Practical sources：实际电压源、实际电流源 $\to$ §6.3、§6.4
\\
p46--p52 & 1.6.3 & Source transformation：五问、公式、2A 演示 $\to$ §7
\\
p53--p58 & 1.6.4 & Example：例题 E1 的五步化简 $\to$ §8
\\
p59--p60 & Review & 复习题、全节小结 $\to$ §9.1、§11
\tabbot
\end{tabular}

这里还有一处必须点明的\textbf{编号分叉}：p3 的课程目录写的是
「1.4 Reference Directions in Electric Circuits；1.5 Ideal Circuit Elements；
1.6 Kirchhoff's Laws」，而 p44 起的页眉写的却是「1.6 Practical sources」。
也就是说，实际电源与电源等效变换被插进了本该是基尔霍夫定律的位置，
基尔霍夫定律整节后移、本册没有讲。复习时按「课件页眉」查页码，按「p3 目录」查知识点归属，
两边对不上是正常的，不是记错。

\subhdr{0.2 记号：先统一一次，后文不再回头}

\par\vspace{2pt}{\small
\noindent\begin{tabular}{@{}p{2.3cm}p{13.4cm}@{}}
\tabtop
\thead{记号} & \thead{约定}
\\
\tabmid
$U,\ I$ & 直流量或恒定量（大写）：$U_S$ 电源电压、$I_S$ 电源电流、$U_{\mathrm{OC}}$ 开路电压
\\
$u,\ i$ & 随时间变化的量（小写）：$u = L\,\mathrm{d}i/\mathrm{d}t$ 里的 $u,i$
\\
$R,\ R_0$ & $R$ 是负载或一般电阻；$R_0$ 专指\textbf{电源内阻}（p14 的原文用 $R_0$，p44 只画不给符号，本讲义统一记 $R_0$）
\\
$V_a,\ U_{ab}$ & $V_a$ 是 a 点\textbf{电位}（相对参考点），$U_{ab}$ 是 a 对 b 的\textbf{电压}，$U_{ab}=V_a-V_b$
\\
参考方向 & 图上画的箭头／正负号，是\textbf{假定}的；真实方向由计算结果的正负反推
\\
关联方向 & 电流参考方向从电压参考极性的「$+$」端\textbf{流入}元件。反之称\textbf{非关联}
\\
$P,\ p$ & 功率。本讲义约定：\textbf{在关联方向下}，$p=ui$ 是元件\textbf{吸收}的功率
\\
$e$ & 电感上的感应电动势，其参考极性与 $u$ 相反（p27 的 $-$ 在上、$+$ 在下），故 $u=-e$
\tabbot
\end{tabular}}

\subhdr{0.3 三处拼接（课件自己的缝）}

\gapitem{编号分叉}{p3 的目录把 1.6 记作 Kirchhoff's Laws，p44 起的页眉却写 1.6 Practical sources。
同一册里「1.6」指两件不同的事。本讲义一律以页眉为准，因为页眉与图、与公式是同一次印刷。}

\gapitem{拼接缺口}{p25 给出的欧姆定律是 $U=RI$，而那张图上\textbf{既没有电流箭头也没有正负极}
（本讲义按 PDF 矢量图元核对：该页只有一个矩形元件与两条引线，无箭头图元）。
紧接着 p27 的三栏表里，三栏的参考方向却全部取成关联方向。于是
「$U=RI$ 什么时候要加负号」这个问题，在 p25 与 p27 之间空了一格，全册没有一页回答。
这是本册最重的一处缺口，§4 专门补它。}

\gapitem{版面跳变}{p39、p40、p41 是「两支理想源并联、电阻并在一旁」的三支并联电路；
p42 把元件值完全相同的 3A、1$\Omega$、3V 换成\textbf{一条无分支的串联回路}
（矢量坐标可证：p39 的 1$\Omega$ 是竖直支路，p42 的 1$\Omega$ 是水平顶边）。
两页之间没有任何一句话说明拓扑变了。§5.3 把两种拓扑分别算到底。}

\callout{\textbf{本讲义与技能骨架的对应关系（§编号不在原位，特此声明）}：\\
lecture-deepdive 的标准骨架是「§5 主例题—§8 公式卡—§9 易错点—§10 缺口—§11 三句话—§12 页码」。
本册课件把「电源模型＋等效变换」整块放在例题之前、又把功率判据夹在中间，
若按骨架硬套，主例题要用到的工具（p44--p52）会出现在例题之后。
因此本讲义把 §5--§7 留给「功率—电源—等效变换」这条工具链，
\textbf{主例题 E1 落在 §8}，§9 收其余例题与「算不下去」的三处，
§10--§12 仍按骨架放缺口清单、三句话、页码对照表。}

\subhdr{0.4 本讲义怎么读}

\callout{
\textbf{第一遍}：只看 §0.1 的表、§8 的 E1 全程、§11 的三句话。这三处能让你把作业做对。\\
\textbf{第二遍}：看 §4（欧姆定律的负号）与 §5（功率的正负）。这两处是全册的枢纽，
考试里只要把电流箭头或电源极性翻一下，答案的符号就全靠它们。\\
\textbf{第三遍}：看 §10 缺口清单。课件哪句话让你卡住，去那一条找页码与补法。}

\subhdr{0.5 版面事实的来源（重要）}

本讲义对「哪个元件接在哪两个节点之间」的判断，不是靠 OCR 猜的，而是按课件的
\textbf{PDF 矢量图元坐标}逐页复原：电阻是带绿色填充的矩形、理想源是红色圆、
导线是宽度 2.28--3.0 的直线段、节点是黑色实心小圆，它们的坐标在同一次印刷里是确定的。
文字一律取自 PDF 文字层。可核对的例子：p39 的 $1\,\Omega$ 是竖矩形（$20\times64$ pt）、
p42 的 $1\,\Omega$ 是横矩形（$61\times20$ pt），故一为并联支路、一为串联顶边；
p19 的开关刀片斜置（断开）、p20 的刀片竖直（闭合）；p44 右栏电压源下方确有内阻矩形。
\textbf{两处仅凭坐标无法判定}的细节已标 [待确认]：
p31 四象限图中三条斜边标注的归属、p32 两张无轴名相图的变量含义（均见 §4.2.2）。
其余各页的连接关系均可按上述规则复算验证。

%=============================================================
\hdr{§1}{电路是什么：功能、组成与三种状态（p2--p16）}

\subhdr{1.1 定义：四个部件与一条闭合路径（p5）}

p5 给的电路定义是这一句：

\par\medskip\noindent
「Circuit: A closed path or loop around which current flows.」（电路：电流流过的闭合路径。）
\par\medskip

页面左侧是实物照片式的电路（两节干电池、一个刀开关、一只灯泡、几段导线），右侧是同一回路的
标准符号图（电源 $E$、开关 S、电阻负载）。课件要你在这两幅图之间认出四件事：
\textbf{电源}（Source，把其他形式的能量变成电能）、\textbf{负载}（Load，把电能变成其他形式）、
\textbf{中间环节}（Switch 与 Wire，控制与连接）、以及\textbf{那是「一条」路径}。

课件只说「闭合」，没说\textbf{为什么必须闭合}。补上这一步，后面 p15 的「开路电流为零」
才不是硬背的结论：

\callout{\textbf{为什么必须闭合}（本讲义补）：在稳态下，电荷既不能凭空产生，也不能在导线某处堆积。
沿回路任取一个横截面，单位时间流过的电荷必须相等，否则该截面就会被越积越多的电荷改变电位，
直到把电流掐断——那就不是稳态了。所以\textbf{沿一条串联回路，电流处处相等}；
一旦有一处断开，整个回路的电流\textbf{同时}为零，而不是「断开处之后才没有电流」。}

\subhdr{1.2 全册标题里的「电路模型」，课件正文一次没出现}

课程标题与文件名都写着「电路模型」，但 p1--p60 没有任何一页出现 model 一词。
这是本册的一处\textbf{形而上缺口}：后面所有公式（$U=IR$、$u=L\,\mathrm{d}i/\mathrm{d}t$、
$P=UI$、电源等效变换）都建立在一个从未被点明的假设上。补一句话：

\callout{\textbf{集总假设（lumped assumption）}（本讲义补，[教材补充]）：
当电路的几何尺寸远小于它工作时电磁波波长的十分之一时，可以把电场效应集中到电容里、
磁场效应集中到电感里、耗散集中到电阻里，于是每个元件只对外露出两个端钮，
整个电路只需要用「端电压 $u$」与「端电流 $i$」两组数描述。
这条假设一旦不成立（高频、长线），$u=L\,\mathrm{d}i/\mathrm{d}t$ 就不再够用，
必须改用分布参数或场的方法。\\
\textbf{为什么必须在这里补}：不承认集总假设，就无法解释「为什么一个元件只要两个端电压与端电流
就能说清它的全部行为」，而这句话正是 p26--p30 三栏表、p33--p36 两种理想源的立论基础。}

\subhdr{1.3 按功能分两类：电力电路与信号电路（p6--p8）}

p6 把电路一分为二，同时提了一个问题：「What is the difference between a power circuit and a
signal circuit?」。p7、p8 分别给出两类电路的功能、特征、成员与关注点：

\par\vspace{2pt}
\noindent\begin{tabular}{@{}p{2.6cm}p{6.6cm}p{6.6cm}@{}}
\tabtop
\thead{对照项} & \thead{电力电路（p7）} & \thead{信号电路（p8）}
\\
\tabmid
Function & Generation, transmission and conversion of electrical energy & Transmission and processing of signals
\\
Characteristics & voltage, current, power ... & voltage, current, power
\\
Members & lighting, power distributing ... & telephone, network, Wi-Fi
\\
Focus & \textbf{efficiency}（效率） & \textbf{effect}，原文括注「Ensure signal transmission quality」
\\
\tabbot
\end{tabular}

读这张表最容易读错的是「Characteristics」一栏：两类电路的物理量\textbf{完全一样}，
都是电压、电流、功率。所以\textbf{区别不在用哪些量，而在电路要完成的任务}：

\warnbox{\textbf{判据（本讲义补）}：在电力电路里，电能是\textbf{产品}，所以关心「送出去多少、损耗多少」；
在信号电路里，电能是\textbf{载体}，所以关心「波形有没有被改变」。
同一只电阻，在电炉里是负载，在分压器里是信号通路的元件——器件不变，任务变了。}

p11 接着问了四个问题，而\textbf{整页只问不答}（这是本册第一处「答案形态」缺口）。逐条答掉：

\par\vspace{2pt}
\noindent\begin{tabular}{@{}p{1.1cm}p{7.1cm}p{7.6cm}@{}}
\tabtop
\thead{问} & \thead{页上原文} & \thead{答案（本讲义补）}
\\
\tabmid
（1）& How many types of circuits based on the purpose and function?
& 两类：电力电路、信号电路（p6）
\\
（2）& What is the difference between a power circuit and a signal circuit?
& 见上表：电能是产品还是载体，关注效率还是效果
\\
（3）& Is the circuit a power circuit or a signal circuit?（配图：电池 $E$ 点一只灯泡）
& \textbf{电力电路}：回路的目的就是把电能变成光与热，灯泡是负载
\\
（4）& The circuit is a simple temperature measuring circuit. Is it a power circuit or a
signal circuit?（配图：电源与一只 mA 表）
& \textbf{信号电路}：虽然也有电源在供电，但回路的目的在于「表针读数」，
电流是信息的载体，电压高低只作为温度的代表
\tabbot
\end{tabular}

\aside{第（4）问是四个问里唯一有争议的一个：mA 表回路确实在耗电，若只看功率就是电力电路。
判据落在「这个回路是为能量还是为信息而存在」上，所以答案是信号电路。
p11 的图只画了「电源＋mA 表」，没有画被测的温度传感器，本答案依据的是题干
「simple temperature measuring circuit」这句话。}

\subhdr{1.4 三个组成部分与「负载增大」到底指什么（p9--p10）}

p9 把电路拆成三块——Sources（电源）、Load（负载）、Intermediate link（中间环节）——
并用一只手电筒做实例：电池是电源、小灯泡是负载、开关与弹簧是中间环节。
p10 复用同一组三块，追加一个问题：「What is the meaning of ``Load Increase''?」，
并在旁边给了答案：

\par\medskip\noindent
「It means the load consumes more electric power, that is, the load current increases.」
\par\medskip

这句话把\textbf{现象}（功率变大）与\textbf{判据}（电流变大）挤在同一行，读起来像循环定义。
它想说的因果链其实有前提，课件没写：

\par\vspace{2pt}{\small
\noindent\begin{tabular}{@{}p{0.9cm}p{14.8cm}@{}}
\tabtop
\thead{步} & \thead{推理（括号内为 p10 图上给定的条件）}
\\
\tabmid
1 & 「负载增大」指的是\textbf{多接入一条并联支路}（p10 的图上 $R_1$、$R_2$ 各 $2\,\Omega$，并在 $U_S$ 两端）
\\
2 & 并联支路增多 $\Rightarrow$ 等效负载电阻\textbf{减小}
\\
3 & 端电压不变（$U_S$ 是理想源）$\Rightarrow$ 电源送出的总电流增大
\\
4 & 功率 $P=UI$ 随电流一起增大
\tabbot
\end{tabular}}

代 p10 图上的数字算一遍（$U_S=10\,\mathrm{V}$，$R_1=R_2=2\,\Omega$）：

\par\vspace{1pt}
\noindent
（1）只接 $R_1$：$I=10/2=5\,\mathrm{A}$，$P=10\times5=50\,\mathrm{W}$；\\
（2）$R_1\parallel R_2$：等效电阻 $1\,\Omega$，$I=10/1=10\,\mathrm{A}$，$P=100\,\mathrm{W}$。
\par\vspace{2pt}

电流翻倍、功率翻倍，与 p10 的文字一致。但\textbf{这两组数在 p34 又出现了一次}
（同一组 $U_S=10\,\mathrm{V}$、$R_1=R_2=2\,\Omega$、$I=5\,\mathrm{A}$ 与 $10\,\mathrm{A}$），
用来解释理想电压源的「电流由外电路决定」。两页合起来才是完整的：
\textbf{「负载增大」是外电路的事，电源只负责把端电压钉住。}

\aside{p10 的推理第 3 步依赖「端电压不变」。若电源有内阻 $R_0$（p14 的模型），
负载增大时端电压会被拉低，电流就不再按 $1/R$ 线性增长——这正是 p44 实际电源要解决的问题。
课件没有在 p10 埋这个伏笔，本讲义把它提前放到这里。}

\warnbox{\textbf{易错}：「负载增大」不是「负载电阻变大」。恰恰相反，负载电阻是变\textbf{小}的
（多并一条支路）。日常说的「用电高峰」也是这个意思：并上去的电器越多，等效电阻越小，总电流越大。}

\subhdr{1.5 基本量清单：p12 只有标题（衔接缺口）}

p12 的页眉已经换成 1.2 Basic quantities，整页只有六行名词：
Current、Electric potential、Voltage、Electromotive force (EMF)、Power、Electric energy。
一页六行、没有定义、没有单位、没有公式。它其实是一张「先自测」的清单，
而本册后面真正用到的只有其中的四个（电流、电压、功率，以及作为电压特例的电动势）。
本讲义按「后面哪里用到」的次序补一列，不重复教科书：

\par\vspace{2pt}
\noindent\begin{tabular}{@{}p{3.4cm}p{3.0cm}p{9.6cm}@{}}
\tabtop
\thead{基本量} & \thead{单位} & \thead{在本册里第一次真用到的地方}
\\
\tabmid
Current $I$ & A & p10 的「负载电流增大」、p13--p16 三态
\\
Electric potential $V$ & V & p17--p20 参考点；$V_a=U_{ad}$ 的定义
\\
Voltage $U$ & V & p22 参考方向三写法、p25 欧姆定律
\\
EMF $E$（$U_S$） & V & p14 的 $U_S$、p15 的 $U_{\mathrm{OC}}=U_S$、p44 的实际源
\\
Power $P$ & W & p37--p43 的吸收／发出判据
\\
Electric energy $W$ & J & 只在 p16 的「过热起火」里隐含出现（$W=Pt$）
\tabbot
\end{tabular}

\subhdr{1.6 三种工作状态：闭合、开路、短路（p13--p16）}

p13 列了三种状态，p14--p16 逐个给出电路图与结论。三页共用同一套元件：
理想电压源 $U_S$、内阻 $R_0$（p14 图下的原文是「$R_0$ is the internal resistance of the source」）、
开关 S、负载 $R$。三态的数据合起来就是\textbf{同一个电源在不同负载下的三个点}：

\par\vspace{2pt}
\noindent\begin{tabular}{@{}p{1.7cm}p{2.9cm}p{3.3cm}p{3.9cm}p{3.2cm}@{}}
\tabtop
\thead{状态} & \thead{等效负载} & \thead{电流} & \thead{端电压} & \thead{电源内部损耗}
\\
\tabmid
闭合 p14 & $R$ 有限 & $I=\dfrac{U_S}{R_0+R}$ & $U=U_S\dfrac{R}{R_0+R}$ & $I^2R_0$
\\
开路 p15 & $R\to\infty$ & $I=0$ & $U=U_{\mathrm{OC}}=U_S$ & $0$
\\
短路 p16 & $R=0$ & $I=I_{\mathrm{SC}}=\dfrac{U_S}{R_0}$ & $U=0$ & $\dfrac{U_S^2}{R_0}$
\tabbot
\end{tabular}

\aside{表中「闭合」一行与「短路」一行的两个式子，课件都没有写：p14 只画图不给式，
p16 只说结论「Short-circuit current is the largest current which may be drawn from the source」，
没有给出 $I_{\mathrm{SC}}=U_S/R_0$。而这条式子正是 p16 那句「最大电流」的定量版本，
也是 p44 实际电源伏安特性与 p48 电源变换公式的来源。本讲义的写法是把 p14 的图与 p16 的结论合成一张表。}

代一组数看清三态的距离（$U_S=10\,\mathrm{V}$，$R_0=0.5\,\Omega$，$R=4.5\,\Omega$）：

\par\vspace{1pt}
\noindent
闭合：$I=\dfrac{10}{0.5+4.5}=2\,\mathrm{A}$，$U=10\times\dfrac{4.5}{5}=9\,\mathrm{V}$，
$P_R=I^2R=4\times4.5=18\,\mathrm{W}$，效率 $\dfrac{18}{10\times2}=90\%$；\\
开路：$I=0$，$U=10\,\mathrm{V}$，$P=0$；\\
短路：$I=\dfrac{10}{0.5}=20\,\mathrm{A}$，$U=0$，$P_{\text{内阻}}=\dfrac{10^2}{0.5}=200\,\mathrm{W}$。
\par\vspace{2pt}

短路时电源内部要烧掉 $200\,\mathrm{W}$，是正常输出 $18\,\mathrm{W}$ 的十一倍——
这就是 p16 那句「may even destroy the source and overheat the connecting wires,
possibly causing a fire」的算术根据。注意它\textbf{与负载无关}：短路电流由电源自己决定。

\warnbox{\textbf{一句话串起三态}：把电流 $I$ 画在横轴、端电压 $U$ 画在纵轴，
三种状态就是\textbf{同一条直线上的三个点}——短路点 $\left(U_S/R_0,\,0\right)$、
工作点 $\left(I,\,U_S-IR_0\right)$、开路点 $\left(0,\,U_S\right)$。
这条直线就是 p44 右侧那条下垂斜线，斜率的负值正是 $R_0$。§6 会把这条直线写成方程。}

%=============================================================
\hdr{§2}{主线一：参考点——电位需要一个海平面（p17--p20）}

\subhdr{2.1 p17 的三个问题为什么不能直接答}

p17 给出的电路是：$10\,\mathrm{V}$ 电源串联 $4\,\Omega$、$3\,\Omega$、$3\,\Omega$ 三个电阻成一条回路，
在 $a$、$b$、$c$ 三个节点上各点了一个红点，问三点的电位（potential）。整页只有三个问句，
没有答案，也没有标电流方向。

\begin{center}
\begin{tikzpicture}[scale=0.82,>=Stealth]
  % 回路
  \draw[cwire] (0,0) -- (0,5.4);
  \draw[cwire] (0,5.4) -- (6.6,5.4);
  \draw[cwire] (6.6,5.4) -- (6.6,0);
  \draw[cwire] (6.6,0) -- (0,0);
  % 电源（左竖边）
  \draw[csrc] (0,2.7) circle (0.34);
  \node[cval] at (0.55,3.25) {$+$};
  \node[cval] at (0.55,2.15) {$-$};
  \node[cval,left] at (-0.15,2.7) {$10\,\mathrm{V}$};
  % 顶部 4Ω
  \draw[cres] (3.0,5.12) rectangle (3.9,5.68);
  \node[cval] at (3.45,6.15) {$4\,\Omega$};
  % 右竖边两个 3Ω
  \draw[cres] (6.32,3.5) rectangle (6.88,4.5);
  \node[cval,right] at (6.98,4.0) {$3\,\Omega$};
  \draw[cres] (6.32,1.0) rectangle (6.88,2.0);
  \node[cval,right] at (6.98,1.5) {$3\,\Omega$};
  % 节点
  \fill (0,5.4) circle (1.8pt) node[above right,clab] {$a$};
  \fill (6.6,5.4) circle (1.8pt) node[right,clab] {$b$};
  \fill (6.6,2.75) circle (1.8pt) node[right,clab] {$c$};
  \fill (0,0) circle (1.8pt) node[below right,clab] {$d$};
  % 接地
  \draw[cwire] (0,0) -- (-0.9,0);
  \draw[cwire] (-1.15,0) -- (-0.65,0);
  \draw[cwire] (-1.08,-0.12) -- (-0.72,-0.12);
  \draw[cwire] (-1.0,-0.24) -- (-0.8,-0.24);
\end{tikzpicture}
\end{center}

第一个问题就是\textbf{凭什么}问「某一点的电位」。一页纸上的三个点，各自接在导线上，
导线里明明有电流在流、电阻上明明有电位降落，为什么还能给每一个点\textbf{一个数}？

\callout{\textbf{电位的定义}：选定一个点作\textbf{参考点}（reference point，惯常取地球的电位为零，
所以也叫「接地」），把其余各点与它之间的电压定义为该点的电位：$V_x = U_{x\,\text{参考点}}$。
换成一句白话：\textbf{电压是两点之间的事，电位是「这一点相对那个被选中的点」的事}。
p17 的图上 $d$ 端画了接地符号，这就是课件的选择——参考点取 $d$，故 $V_d = 0$。}

\subhdr{2.2 p18 的三个数是怎么来的}

p18 直接给出三个结果，一个中间步骤也没有（这是本册第二处「答案形态」缺口）：

\par\medskip\noindent
$V_a = U_{ad} = 10\,\mathrm{V}$，\quad $V_b = U_{bd} = 6\,\mathrm{V}$，\quad
$V_c = U_{cd} = 3\,\mathrm{V}$。
\par\medskip

三式都写成 $V_x = U_{xd}$ 的形式，这是在提醒你：\textbf{电位就是「对参考点的电压」}。
至于 $10$、$6$、$3$ 这三个数，来自两件事，一件用不着算，一件必须算：

\par\vspace{2pt}
\noindent\begin{tabular}{@{}p{0.9cm}p{6.2cm}p{8.4cm}@{}}
\tabtop
\thead{步} & \thead{做法} & \thead{说明}
\\
\tabmid
1 & 直接读出 $V_a$ & $a$ 与 $d$ 之间只隔着一个理想电压源，理想源的端电压不随电流变，
故 $U_{ad} = 10\,\mathrm{V}$，与后面接了什么电阻无关
\\
2 & 算回路电流 & 三个电阻串联：$\sum R = 4+3+3 = 10\,\Omega$，
$I = 10/10 = 1\,\mathrm{A}$（稳态串联回路电流处处相等，§1.1）
\\
3 & 逐段算压降 & $U_{ab} = IR = 1\times4 = 4\,\mathrm{V}$；
$U_{bc} = 1\times3 = 3\,\mathrm{V}$；$U_{cd} = 1\times3 = 3\,\mathrm{V}$
\\
4 & 沿途累减 & $V_b = V_a - U_{ab} = 10-4 = 6\,\mathrm{V}$；
$V_c = V_b - U_{bc} = 6-3 = 3\,\mathrm{V}$；$V_d = V_c - U_{cd} = 3-3 = 0\,\mathrm{V}$
\tabbot
\end{tabular}

最后一步回到 $V_d = 0$，与「$d$ 是参考点」自洽，这是最好用的一种自查：
\textbf{沿回路绕一圈，所有压降之和必须等于起点与终点之间已知的电位差}。

\aside{第 4 步用的是「沿电流方向走，电位下降」这条约定。它与 p22 的电压参考方向一致：
电压箭头的箭头端就是电位较低的一端。若把电流方向取反（比如按 §3 的惯例改成从 $d$ 流向 $a$），
则 $I$ 变成 $-1\,\mathrm{A}$，$U_{ab}$ 变成 $-4\,\mathrm{V}$，累减仍然是「减一个负数」，
答案不变。这一点在 §3.3 会讲清。}

\subhdr{2.3 换参考点：所有电位一起平移（本讲义补）}

课件给的是「$d$ 作参考点」这一种选择，但真实电路里换参考点是家常便饭
（示波器的地夹子夹在哪，哪就是零）。换参考点之后各点电位怎么变？p18 没有回答，
而这是本章唯一一条能写成定理的结论：

\subhdr{定理（参考点平移）}

设参考点由 $d$ 改为 $e$。则对任意一点 $x$，新的电位 $V'_x$ 满足

\[ V'_x = V_x + U_{ed}, \]

即\textbf{所有点的电位同时加上同一个常数} $U_{ed}$；而任意两点之间的电压不变：
$V'_x - V'_y = V_x - V_y$。

\textbf{证明}：按定义，$V_x = \varphi_x - \varphi_d$，$V'_x = \varphi_x - \varphi_e$，
两式相减得

\[ V'_x - V_x = \varphi_e - \varphi_d = U_{ed}, \]

与 $x$ 无关，故是一次整体平移。两点的电压相减时这个常数自动抵消，所以电压不变。$\square$

\aside{本定理是本讲义补的，p18 只给三个数、不谈换参考点。它之所以值得写成定理，
是因为「电位可以整体平移而电压不能」这条性质，正是 §2.4 区分电位与电压的唯一依据。}

把它代到 p18 的数上：若把参考点从 $d$ 改到 $c$，则 $U_{cd} = 0 - 3 = -3\,\mathrm{V}$，每点加 $-3\,\mathrm{V}$：

\par\vspace{2pt}
\noindent\begin{tabular}{@{}p{2.2cm}p{3.4cm}p{3.4cm}p{4.4cm}@{}}
\tabtop
\thead{点} & \thead{$d$ 为参考（p18）} & \thead{$c$ 为参考（本讲义）} & \thead{相邻两点电压（两列相同）}
\\
\tabmid
$a$ & $10\,\mathrm{V}$ & $7\,\mathrm{V}$ & $U_{ab} = 4\,\mathrm{V}$
\\
$b$ & $6\,\mathrm{V}$ & $3\,\mathrm{V}$ & $U_{bc} = 3\,\mathrm{V}$
\\
$c$ & $3\,\mathrm{V}$ & $0\,\mathrm{V}$ & $U_{cd} = 3\,\mathrm{V}$
\\
$d$ & $0\,\mathrm{V}$ & $-3\,\mathrm{V}$ & ---
\tabbot
\end{tabular}

\warnbox{\textbf{由此得到本章最该记住的一句}：\textbf{电位会变，电压不变}。
题目里出现「求某点电位」时，第一件事是找参考点；题目问「求某两点之间的电压」时，
参考点取在哪都不影响答案。p17 的问法是前者，所以参考点必须先定；
p19、p20 的问法也是前者（求 $V_A$），参考点由接地点给出。}

\subhdr{2.4 电位与电压的分界线：双下标的反对称性}

有了 §2.3 的定理，就能把两个概念彻底分开，并顺带补上 p22 没有写的一条恒等式：

\par\medskip\noindent
$V_x$：依附于参考点，换参考点要整体平移；\quad
$U_{xy} = V_x - V_y$：与参考点无关，是两点的\textbf{差}。
\par\medskip

于是立刻有

\[ U_{ab} = V_a - V_b = -(V_b - V_a) = -U_{ba} . \]

这条 $U_{ab} = -U_{ba}$ 是双下标写法的全部技术内容：\textbf{把下标的两个字母对调，就等于加一个负号}。
p22 只画了「$U_{AB}$ 的方向由 A 到 B」，没有写这条反对称性，
而后面 p23、p24 的电路里 $U_{ab}$ 与 $U_{ba}$ 会同时出现，不记住这条就要乱。代 §2.2 的数：$U_{ab} = 10-6 = 4\,\mathrm{V}$，
$U_{ba} = 6-10 = -4\,\mathrm{V}$，确为相反数。

\subhdr{2.5 p19 与 p20：两问都没有答案，全程补上}

p19、p20 是同一张图的两问：给定 $R_1=6\,\mathrm{k}\Omega$、$R_2=4\,\mathrm{k}\Omega$、
$R_3=20\,\mathrm{k}\Omega$，求 $A$ 点电位，先是「when S is open」，后是「when S is closed」。
\textbf{课件两页都没有给出答案}，也没有任何演算。

先把拓扑复原清楚。按该页矢量图元的坐标：$R_1$ 是左上方的竖直电阻、$R_2$ 是中部的水平电阻、
$R_3$ 是右下方的竖直电阻；上端钮标 $-12\,\mathrm{V}$、下端钮标 $+12\,\mathrm{V}$；
$R_1$ 与 $R_2$ 的交点（记作 $M$）向下经开关 S 接到接地符号。

\begin{center}
\begin{tikzpicture}[scale=0.82,>=Stealth]
  % 端钮
  \draw[csrc,red] (2.6,6.4) circle (0.26);
  \node[cval,left] at (2.2,6.4) {$-12\,\mathrm{V}$};
  \draw[csrc,red] (7.6,0.5) circle (0.26);
  \node[cval,below] at (7.6,0.1) {$+12\,\mathrm{V}$};
  % R1 竖直
  \draw[cwire] (2.6,6.4) -- (2.6,5.6);
  \draw[cres] (2.32,4.6) rectangle (2.88,5.6);
  \node[cval,left] at (2.2,5.1) {$R_1$};
  \draw[cwire] (2.6,4.6) -- (2.6,3.9);
  \fill (2.6,3.9) circle (1.8pt);
  \node[cval,below left] at (2.55,3.85) {$M$};
  % R2 水平
  \draw[cwire] (2.6,3.9) -- (3.6,3.9);
  \draw[cres] (3.6,3.62) rectangle (4.9,4.18);
  \node[cval,above] at (4.25,4.3) {$R_2$};
  \draw[cwire] (4.9,3.9) -- (7.6,3.9);
  % A 点与 R3
  \fill (7.6,3.9) circle (1.8pt) node[right,clab] {$A$};
  \draw[cwire] (7.6,3.9) -- (7.6,3.1);
  \draw[cres] (7.32,1.5) rectangle (7.88,3.1);
  \node[cval,right] at (7.98,2.3) {$R_3$};
  \draw[cwire] (7.6,1.5) -- (7.6,0.5);
  % 开关 S 与地
  \draw[cwire] (2.6,3.9) -- (2.6,2.4);
  \draw[thick,red] (2.6,2.4) -- (1.5,1.75);
  \fill (2.6,2.4) circle (1.2pt);
  \draw[cwire] (2.6,1.35) -- (2.6,0.9);
  \draw[cwire] (2.1,0.9) -- (3.1,0.9);
  \draw[cwire] (2.2,0.75) -- (3.0,0.75);
  \draw[cwire] (2.3,0.6) -- (2.9,0.6);
  \node[cval,left] at (2.3,1.95) {S};
\end{tikzpicture}
\end{center}

\subhdr{第一步：认参考点}

图上有接地符号，故接地点的电位为 $0\,\mathrm{V}$。这一点比「$+12\,\mathrm{V}$ 与 $-12\,\mathrm{V}$」更重要：
$-12\,\mathrm{V}$ 与 $+12\,\mathrm{V}$ 是\textbf{两个端钮对地的电位}，不是两个「电源」，
不要把它们读成「负电源往下掉」。

\subhdr{第二步：开关断开时（p19）}

S 断开 $\Rightarrow$ $M$ 点不与地相连 $\Rightarrow$ 从 $-12\,\mathrm{V}$ 到 $+12\,\mathrm{V}$ 只剩一条串联链：
$R_1 - R_2 - R_3$，共 $30\,\mathrm{k}\Omega$，两端电位差 $24\,\mathrm{V}$。

\[ I = \frac{24}{30} = 0.8\,\mathrm{mA} \quad(\text{由 } +12\,\mathrm{V} \text{ 端流向 } -12\,\mathrm{V} \text{ 端}) \]

用 $+12\,\mathrm{V}$ 端作起点、沿 $R_3$ 上溯到 $A$（这一步是顺着电流走，电位下降）：

\[ V_A = 12 - 0.8\times20 = 12-16 = -4\,\mathrm{V}. \]

\textbf{自查}：继续往上过 $R_2$ 得 $V_M = -4 - 0.8\times4 = -7.2\,\mathrm{V}$；
再过 $R_1$ 得 $-7.2-0.8\times6 = -12\,\mathrm{V}$，正好回到上端钮。$\checkmark$

\subhdr{第三步：开关闭合时（p20）}

S 闭合 $\Rightarrow$ $M$ 被钳到 $0\,\mathrm{V}$。此时电路裂成两半：
左边的 $R_1$ 跨在 $-12\,\mathrm{V}$ 与地之间，右边的 $A$ 点只与 $R_2$（接 $M=0$）和 $R_3$（接 $+12\,\mathrm{V}$）有关。

于是 $A$ 点的问题变成一个\textbf{分压器}：$0\,\mathrm{V}$ 与 $+12\,\mathrm{V}$ 之间串着 $R_2=4\,\mathrm{k}\Omega$
与 $R_3=20\,\mathrm{k}\Omega$，$A$ 是分压点，

\[ V_A = 12\times\frac{R_2}{R_2+R_3} = 12\times\frac{4}{24} = 2\,\mathrm{V}. \]

\textbf{自查}：流过 $R_3$（由上往下）与 $R_2$（由 $A$ 往 $M$）的是同一个电流
$12/24 = 0.5\,\mathrm{mA}$，故 $V_M + 0.5\times4 = 0+2 = V_A$ $\checkmark$。
左边的 $R_1$ 上另有 $12/6 = 2\,\mathrm{mA}$ 由地流向 $-12\,\mathrm{V}$ 端钮，\textbf{但不进入 $A$ 点}，
所以对 $V_A$ 没有影响。

\par\vspace{2pt}
\noindent\begin{tabular}{@{}p{2.4cm}p{3.6cm}p{3.6cm}p{4.4cm}@{}}
\tabtop
\thead{情况} & \thead{串联总电阻} & \thead{电流} & \thead{$V_A$}
\\
\tabmid
S 断开（p19） & $6+4+20=30\,\mathrm{k}\Omega$ & $0.8\,\mathrm{mA}$ & $-4\,\mathrm{V}$
\\
S 闭合（p20） & 见上（$R_2+R_3=24\,\mathrm{k}\Omega$） & $0.5\,\mathrm{mA}$（$A$ 支路） & $+2\,\mathrm{V}$
\tabbot
\end{tabular}

\warnbox{\textbf{本题的两个坑}（都指回具体一步）：\\
（1）\textbf{把「$-12\,\mathrm{V}$」当成「接在参考点上」}。它离参考点有 $12\,\mathrm{V}$ 的差距，
不是 $0$。凡是端钮旁边直接写电压数值的，都是「该端钮对地的电位」，要当成一个理想的固定电位。\\
（2）\textbf{合上开关后仍然去算那条串联链}。S 一合，$M$ 被钳到 $0$，原来的「一条 $30\,\mathrm{k}\Omega$ 串联链」
就不存在了，$R_1$ 与右边两支路已经解耦。判断串并联永远是\textbf{先看节点被钳在哪里}。}

\aside{$V_A$ 从 $-4\,\mathrm{V}$ 跳到 $+2\,\mathrm{V}$，靠的是「接地把 $M$ 钳到 $0$」这一件事。
若把接地点换到 $A$（题目就变成求 $V_M$），按 §2.3 的平移定理，
全部电位同时减去 $2\,\mathrm{V}$：$V_A = 0$、$V_M = -2\,\mathrm{V}$。可见\textbf{开关的作用是改变拓扑，
接地点的作用只是平移读数}，两件事不能混。}

%=============================================================
\hdr{§3}{主线二：参考方向——为什么方向可以「随便假设」（p21--p24）}

\subhdr{3.1 全册最反直觉的一步：先假一个方向，再算}

p21 只写了两行：「（1）The directions we concern； （2）What is reference direction? assumed」。
第二行的 assumed（假定）一词是全册最容易被略过、又最要紧的一个词。
它意味着图上的箭头\textbf{不是}电流的真实方向，而是解方程之前\textbf{自己先选定}的一个方向。

p23 的 Example 1 说：「If we know all the real directions, just use the real directions as the
reference directions.」——知道真实方向时，就把真实方向取作参考方向（省事）。
p24 的 Example 2 说：不知道时怎么办——

\par\medskip\noindent
「(1) Before solving the problem, assume a current direction as a reference direction.
(2) Do the calculation. (3) If the calculation result $I$ is positive, the actual direction is
consistent with the reference direction; If the calculation result $I$ is negative, the actual
direction is opposite to the reference direction.」
\par\medskip

这三步里，第 (1) 步和第 (2) 步都是常规操作，\textbf{全部意义在第 (3) 步}：结果的正负号不是「算错了」，
而是「方向报告」。课件把三步并列写出，容易让人以为第 (3) 步是个附加说明；其实它是这套约定的\textbf{兑现方式}。

\subhdr{3.2 三种写法与它们之间的换算（p22）}

同一个电压参考方向，p22 给了三种写法，并提醒「Notice」：

\par\vspace{2pt}
\noindent\begin{tabular}{@{}p{1.0cm}p{3.6cm}p{5.4cm}p{5.6cm}@{}}
\tabtop
\thead{序} & \thead{写法} & \thead{图上怎么画} & \thead{译成一句话}
\\
\tabmid
（1） & by arrows & 元件上方一支箭头，旁边标 $U$ & 「箭头由 $+$ 指向 $-$」，
即\textbf{箭尾端为正、箭头端为负}
\\
（2） & by positive and negative polarity & 元件两端标 $+$ 与 $-$ & 「$+$ 端电位高于 $-$ 端」
\\
（3） & by double subscripts & 两端标 $A$、$B$，标 $U_{AB}$ & 「$U_{AB}=V_A-V_B$」，
下标\textbf{先正后负}
\tabbot
\end{tabular}

\begin{center}
\begin{tikzpicture}[scale=0.8,>=Stealth]
  % (1) 箭头
  \draw[cwire] (0,1) -- (1.9,1);
  \draw[cres] (1.9,0.7) rectangle (3.1,1.3);
  \draw[cwire] (3.1,1) -- (5,1);
  \fill[black] (1.9,1) circle (1.4pt); \fill[black] (3.1,1) circle (1.4pt);
  \draw[carr,red] (2.1,1.85) -- (2.9,1.85);
  \node[cval,red] at (2.5,2.2) {$U$};
  \node[cval] at (2.5,-0.35) {（1）箭头};
  % (2) 正负极
  \draw[cwire] (6.4,1) -- (8.3,1);
  \draw[cres] (8.3,0.7) rectangle (9.5,1.3);
  \draw[cwire] (9.5,1) -- (11.4,1);
  \fill[black] (8.3,1) circle (1.4pt); \fill[black] (9.5,1) circle (1.4pt);
  \node[cval,red] at (7.9,1.55) {$+$};
  \node[cval,red] at (9.9,1.55) {$-$};
  \node[cval,red] at (8.9,2.0) {$U$};
  \node[cval] at (8.9,-0.35) {（2）正负极};
  % (3) 双下标
  \draw[cwire] (12.8,1) -- (14.7,1);
  \draw[cres] (14.7,0.7) rectangle (15.9,1.3);
  \draw[cwire] (15.9,1) -- (17.8,1);
  \fill[black] (14.7,1) circle (1.4pt); \fill[black] (15.9,1) circle (1.4pt);
  \node[cval] at (14.3,1.6) {$A$};
  \node[cval] at (16.3,1.6) {$B$};
  \node[cval] at (15.3,2.0) {$U_{AB}$};
  \node[cval] at (15.3,-0.35) {（3）双下标};
\end{tikzpicture}
\end{center}

三种写法是\textbf{同一件事的三种语言}，可以机械互译（p22 只把三种写法并排放，没有给换算规则，
这是本节的一处「机制缺口」）：

\callout{\textbf{互译规则（本讲义补）}：
（1）箭头 $\to$（2）正负极：箭尾标 $+$、箭头标 $-$；
（2）正负极 $\to$（3）双下标：\textbf{先 $+$ 后 $-$}，即 $U_{+\,-}$；
（3）双下标 $\to$（1）箭头：箭头由第一个下标指向第二个下标。
再补一条三种写法都适用的恒等式：$U_{AB} = -U_{BA}$（§2.4）。
\textbf{一页之内只能选一种写法定下来}，中途换语言而不换符号，是本章最常见的自伤。}

\subhdr{3.3 为什么「假设」不会算错：线性与坐标系}

学生最常问的一句是「我凭什么知道该往哪边画箭头？」答案分两层：

\textbf{第一层（记号层）}：参考方向不是物理量，是\textbf{坐标系}。就像在数轴上规定向右为正之后，
「走了 3 米向左」记作 $-3$ 米——没有人会说这个记录「错了」。箭头画在哪一边，只是规定了
「正号代表哪个方向」。

\textbf{第二层（数学层）}：电路方程对 $I$ 是\textbf{一次}的（欧姆定律 $U=RI$、基尔霍夫定律的求和都是线性的），
把 $I$ 换成 $-I$、同时把所有相关的电压极性一起翻转，方程在形式上完全不变。
所以假设错方向，解出来的只是「同一个物理状态的另一套坐标读数」，数字差一个负号，物理不变。

\warnbox{\textbf{但有一个代价必须付}：负号的\textbf{解释}。算出 $I=-2\,\mathrm{A}$ 时，不能写「电流是负的」了事，
要写成「电流 $2\,\mathrm{A}$，方向与图上箭头\textbf{相反}」。同理，$P=-9\,\mathrm{W}$ 要说成「发出 $9\,\mathrm{W}$」
（§5.1）。\textbf{参考方向给了你自由，代价是每次读答案都要翻译一次}。}

\subhdr{3.4 一个两行的练习：把三种写法来回翻译}

设元件两端 $A$、$B$，已知 $V_A = 6\,\mathrm{V}$、$V_B = 2\,\mathrm{V}$。

\par\vspace{1pt}
\noindent
（1）双下标写法：$U_{AB} = V_A - V_B = 4\,\mathrm{V}$；\\
（2）极性写法：$A$ 端标 $+$、$B$ 端标 $-$，旁边写 $U = 4\,\mathrm{V}$；\\
（3）箭头写法：箭头由 $A$ 指向 $B$（箭尾在 $+$ 端），旁标 $U = 4\,\mathrm{V}$；\\
（4）若改成问 $U_{BA}$：$U_{BA} = V_B - V_A = -4\,\mathrm{V}$，箭头方向也要跟着反过来。
\par\vspace{2pt}

四种写法说的是同一件事。\textbf{唯一会出错的情形是「改了一半」}：把 $U_{AB}$ 写成 $U_{BA}$
却忘了加负号，或把箭头翻过来却仍然按原来的极性读数值。

%=============================================================
\hdr{§4}{两条线合流：关联与非关联，欧姆定律的负号从哪来（p25--p30）}

\subhdr{4.1 p25 是一张「填空题」，不是一条定律}

p25 的版面只有三样东西：标题「How to determine the reference direction?」、
「Example 3:」、以及一个二端元件的符号和一行公式 $U=RI$。
按该页矢量图元核对，\textbf{图上没有电流箭头，也没有正负极性标记}——它只是一段带两个端钮的元件。

这张「空图」问的正是：$U=RI$ 要成立，图上缺了哪两笔？答案是缺两个参考方向，
而且它们必须\textbf{配套}。这就是本册最重的一处缺口，先补定义：

\callout{\textbf{定义（关联参考方向）}：对一个二端元件，若电流的参考方向是
\textbf{从电压参考极性的「$+$」端流入元件、从「$-$」端流出}，则称 $u$ 与 $i$ 取\textbf{关联参考方向}
（associated reference directions）；若电流是从「$+$」端\textbf{流出}，则称\textbf{非关联}。
注意这里判断的始终是\textbf{参考方向}，与真实电流往哪边流无关。}

\subhdr{定理（欧姆定律的符号形式）}

设线性电阻的阻值为 $R$（$R>0$），其两端电压为 $u$、电流为 $i$。若 $u$、$i$ 取关联参考方向，则

\[ u = Ri ; \]

若取非关联参考方向，则

\[ u = -Ri . \]

\textbf{证明}：把元件的两个端钮记作 $a$（取电压参考极性的 $+$ 端）与 $b$（$-$ 端），故
$u = \varphi_a - \varphi_b$。先设电流的参考方向由 $a$ 经元件流向 $b$（即关联）。
这时 $i>0$ 的含义就是「正电荷确实由 $a$ 流向 $b$」，而电荷沿导体流动时电位单调下降，
故 $a$ 的电位高于 $b$，且落差正比于流过的电流与电阻：$\varphi_a - \varphi_b = Ri$，即 $u = Ri$。

再设电流的参考方向取成由 $b$ 流向 $a$（非关联），并把这个记号下的读数写作 $i'$。
同一个物理电流在两种记法下只差一个符号：$i' = -i$。把它代入上一段的结论：

\[ u = Ri = R(-i') = -Ri' . \qquad\square \]

\aside{这个证明只用了一件事：\textbf{电荷顺着电流方向走时电位下降}。它与 §2.2 用的「沿电流方向走电位下降」
是同一条约定，也是 p22 的「电压箭头由 $+$ 指向 $-$」的来历。换句话说，
§2 定下的记号约定，在 §4 兑现成了符号规则。}

\subhdr{数值小例：两条路必须给出同一个 $u$}

取 $R = 2\,\Omega$，已知真实电流 $3\,\mathrm{A}$ 由 $a$ 流向 $b$，$a$ 为 $+$ 端。

\par\vspace{2pt}
\noindent\begin{tabular}{@{}p{5.0cm}p{4.2cm}p{5.4cm}@{}}
\tabtop
\thead{记法} & \thead{读数} & \thead{代入哪一式}
\\
\tabmid
电流箭头取 $a\to b$，属关联 & $i = 3\,\mathrm{A}$ & $u = Ri = 2\times3 = 6\,\mathrm{V}$
\\
电流箭头取 $b\to a$，属非关联 & $i = -3\,\mathrm{A}$ & $u = -Ri = -2\times(-3) = 6\,\mathrm{V}$
\tabbot
\end{tabular}

两条路都得到 $u = 6\,\mathrm{V} = \varphi_a-\varphi_b$，物理一致。这个 20 秒的小例值得反复看：
\textbf{负号不是新物理，它是「记号与物理方向不一致」的记账费}。
把它记成一句话：\textbf{电流进 $+$ 端，$u=Ri$；电流出 $+$ 端，$u=-Ri$。}

\begin{center}
\begin{tikzpicture}[scale=0.8,>=Stealth]
  % 左：关联
  \draw[cwire] (0,2) -- (1.4,2);
  \draw[cres] (1.4,1.6) rectangle (3.2,2.4);
  \draw[cwire] (3.2,2) -- (4.6,2);
  \node[cval,red] at (1.25,2.55) {$+$};
  \node[cval,red] at (3.35,2.55) {$-$};
  \node[cval,red] at (2.3,3.1) {$u$};
  \draw[carr,blue] (0.2,1.15) -- (1.5,1.15);
  \node[cval,blue] at (0.85,0.65) {$i$};
  \fill[black] (1.4,2) circle (1.3pt); \fill[black] (3.2,2) circle (1.3pt);
  \node[cval] at (2.3,-0.05) {关联：$u=Ri$};
  % 右：非关联
  \draw[cwire] (7.0,2) -- (8.4,2);
  \draw[cres] (8.4,1.6) rectangle (10.2,2.4);
  \draw[cwire] (10.2,2) -- (11.6,2);
  \node[cval,red] at (8.25,2.55) {$+$};
  \node[cval,red] at (10.35,2.55) {$-$};
  \node[cval,red] at (9.3,3.1) {$u$};
  \draw[carr,blue] (10.9,1.15) -- (7.7,1.15);
  \node[cval,blue] at (9.3,0.65) {$i$};
  \fill[black] (8.4,2) circle (1.3pt); \fill[black] (10.2,2) circle (1.3pt);
  \node[cval] at (9.3,-0.05) {非关联：$u=-Ri$};
\end{tikzpicture}
\end{center}

\subhdr{4.2 课件 p27--p30 的三栏表，逐栏读出参考方向}

p27 起是同一张表连着四页，每页只换底部那句结论：三栏分别是 Resistor、Inductor、Capacitor，
每栏一个符号图 + 一条伏安关系。这四页是全册密度最高的地方，也是负号最容易读丢的地方。
按矢量图元逐栏判读：

\par\vspace{2pt}
\noindent\begin{tabular}{@{}p{1.5cm}p{4.0cm}p{3.9cm}p{3.0cm}p{2.4cm}@{}}
\tabtop
\thead{栏} & \thead{图上标注} & \thead{是否关联} & \thead{课件公式} & \thead{同义写法}
\\
\tabmid
R & $U$ 箭头向下（$+$ 在上）；$I$ 沿顶边向右、由元件顶端流入 & 关联 &
$U = IR$ & $R = U/I$
\\
L & $u$ 的 $+$ 在上；$i$ 由顶端流入 & 关联 &
$u = -e = L\dfrac{\mathrm{d}i}{\mathrm{d}t}$ & $e = -L\dfrac{\mathrm{d}i}{\mathrm{d}t}$
\\
C & $u$ 的 $+$ 在上；$i$ 由上方流入上极板 & 关联 &
$i = \dfrac{\mathrm{d}q}{\mathrm{d}t} = C\dfrac{\mathrm{d}u}{\mathrm{d}t}$ & $q = Cu$
\tabbot
\end{tabular}

\begin{center}
\begin{tikzpicture}[scale=0.78,>=Stealth]
  % 电阻栏：U 向下箭头、I 向右
  \draw[cwire] (0.6,2.6) -- (3.0,3.5);
  \draw[cwire] (0.6,1.0) -- (3.0,2.1);
  \draw[cwire] (3.0,2.1) -- (3.0,3.5);
  \draw[cres] (2.72,2.42) rectangle (3.28,3.18);
  \draw[carr,blue] (1.5,3.05) -- (2.4,3.36);
  \node[cval,blue] at (1.95,3.55) {$i$};
  \draw[carr,red] (1.15,2.45) -- (1.15,1.45);
  \node[cval,red] at (0.65,1.95) {$u$};
  \node[cval] at (1.9,0.25) {R：$u=Ri$};
  % 电感栏
  \draw[cwire] (5.4,3.5) -- (7.4,3.5);
  \draw[cwire] (7.4,3.5) -- (7.4,3.0);
  \draw[decorate,decoration={coil,segment length=1.1mm,amplitude=1.3mm},thick]
       (7.4,3.0) -- (7.4,1.6);
  \draw[cwire] (7.4,1.6) -- (7.4,1.0);
  \draw[cwire] (5.4,1.0) -- (7.4,1.0);
  \draw[carr,blue] (6.1,3.5) -- (6.9,3.5);
  \node[cval,blue] at (6.5,3.95) {$i$};
  \node[cval,red] at (6.85,3.05) {$+$};
  \node[cval,red] at (6.85,1.25) {$-$};
  \node[cval,red] at (6.55,2.3) {$u$};
  \node[cval,green!45!black] at (8.0,2.95) {$-$};
  \node[cval,green!45!black] at (8.0,1.35) {$+$};
  \node[cval,green!45!black] at (8.35,2.3) {$e$};
  \node[cval] at (6.4,0.25) {L：$u=L\,\mathrm{d}i/\mathrm{d}t=-e$};
  % 电容栏
  \draw[cwire] (9.8,3.5) -- (11.4,3.5);
  \draw[cwire] (11.4,3.5) -- (11.4,2.75);
  \draw[thick,red!70!black] (10.9,2.75) -- (11.9,2.75);
  \draw[thick,red!70!black] (10.9,2.45) -- (11.9,2.45);
  \draw[cwire] (11.4,2.45) -- (11.4,1.0);
  \draw[cwire] (9.8,1.0) -- (11.4,1.0);
  \draw[carr,blue] (11.4,3.3) -- (11.4,2.9);
  \node[cval,blue] at (11.85,3.15) {$i$};
  \node[cval,red] at (11.15,2.98) {$+$};
  \node[cval,red] at (11.15,2.2) {$-$};
  \node[cval,red] at (10.6,2.6) {$u$};
  \node[cval] at (10.8,0.25) {C：$i=C\,\mathrm{d}u/\mathrm{d}t$};
\end{tikzpicture}
\end{center}

三栏全部取的是\textbf{关联}方向，所以三式都不带负号——这是它们能写成课件那种形式的唯一原因。
把任何一栏的电流箭头翻过来，该栏的公式就要在前面加一个负号。课件没有一页写这句话。

\subhdr{4.2.1 电感那一栏的两个负号，分别是什么}

$u=-e=L\,\mathrm{d}i/\mathrm{d}t$ 里有两条信息，课件压成了一行：

\par\vspace{2pt}
\noindent\begin{tabular}{@{}p{0.9cm}p{13.6cm}@{}}
\tabtop
\thead{步} & \thead{内容}
\\
\tabmid
1 & \textbf{物理定律}：电感的端电压与电流的\textbf{变化率}成正比，比例系数是电感 $L$。
关联方向下写成 $u = L\,\mathrm{d}i/\mathrm{d}t$。电感与电阻的区别全在这里：
电阻看「电流有多大」，电感看「电流变得有多快」。恒定电流下电感相当于导线（$u=0$）。
\\
2 & \textbf{记号关系}：$e$ 是感应电动势。按「电动势」的惯例，它的参考极性画成
\textbf{上 $-$、下 $+$}，与 $u$ 的「上 $+$、下 $-$」正好相反，故 $e=-u$，
于是 $u=-e$。这一行里没有新物理，它只是两种惯例的汇率。
\\
\tabbot
\end{tabular}

\aside{为什么物理上要把 $e$ 画反？因为楞次定律：感应电动势总是\textbf{反抗}电流的变化。
按「电动势从 $-$ 端指向 $+$ 端」的惯例画，反抗就意味着它的极性与「电流流入端为正」的电压极性相反。
若把 $e$ 直接画成与 $u$ 同极性，那就该写 $e = L\,\mathrm{d}i/\mathrm{d}t$，
「负号」也随之消失——所以那一行的负号，本质上是\textbf{惯例的负号}，不是新定律的负号。}

\subhdr{4.2.2 p31 的第四个元件：忆阻器补的是哪一格}

p27--p30 的三栏表只讲了三个无源元件，p31 补上第四个。该页的做法是把四个基本量
——电压 $v$、电流 $i$、电荷 $q$、磁通 $\phi$——摆成十字形的四个端点，每一对端点之间放一个元件，
它们的关系都是\textbf{微分形式}的：

\par\vspace{2pt}
\noindent\begin{tabular}{@{}p{3.2cm}p{4.6cm}p{7.4cm}@{}}
\tabtop
\thead{元件} & \thead{端口关系} & \thead{读法}
\\
\tabmid
Resistor & $\mathrm{d}v = R\,\mathrm{d}i$ & 电压与电流之间只差一个常数
\\
Capacitor & $\mathrm{d}q = C\,\mathrm{d}v$ & 电荷随电压累积
\\
Inductor & $\mathrm{d}\phi = L\,\mathrm{d}i$ & 磁通随电流累积
\\
Memristor & $\mathrm{d}\phi = M\,\mathrm{d}q$ & \textbf{磁通与电荷之间}的那个「缺口」
\tabbot
\end{tabular}

前三个都是中学就熟的，第四个是忆阻器：它把「磁通—电荷」这一对补上，四种元件才两两配齐。
p31 同时给了三条史实：「When: 1971」「Who: Leon Chua」「Practical: 2008 HP Lab」，
而「Characteristics:」与「Applications:」两条\textbf{在页面上是空的}（只有冒号）。
这是本册一处不必算就能看出的缺口，按 [教材补充] 补两行：

\callout{\textbf{[教材补充] 忆阻器的特性与应用}（p31 的（4）（5）两栏课件留空）：\\
\textbf{（4）特性}：忆阻值 $M = \mathrm{d}\phi/\mathrm{d}q$，量纲是「磁通／电荷」，即
V$\cdot$s 除以 A$\cdot$s $=$ V／A，\textbf{正是欧姆}。它与电阻的本质区别是 $M$ 不是常数，
而是流过它的电荷总量的函数，所以\textbf{断电后 $M$ 保持原值}——这就是「记忆」二字的来历。\\
\textbf{（5）应用}：非易失存储（RRAM）、模拟突触的神经形态计算，以及把线性电路理论推广到非线性
（p32 那两张缠成「$\infty$」形的相图就是这个方向的插图）。p31 的版面已经点到
「Integrated Circuit : Graphene」与「Chaos: From linear science to Non-linear science」两个方向。}

\aside{三个无源元件的公式（$u=Ri$、$u=L\,\mathrm{d}i/\mathrm{d}t$、$i=C\,\mathrm{d}u/\mathrm{d}t$）
与忆阻器的 $\mathrm{d}\phi=M\,\mathrm{d}q$ 是同一套路的四种特例：\textbf{都是「某一对量的微分之比」}。
p31 的十字图想说的就是这句话，只是课件没有把这句总结写出来，只摆了图。}

\gapitem{[待确认]（p31）}{四象限图中三条斜边上还各有一处标注（形如 $\mathrm{d}v=R\,\mathrm{d}i$、
$\mathrm{d}\phi=L\,\mathrm{d}i$、$\mathrm{d}\phi=M\,\mathrm{d}q$）。该图由多条斜线与多个文本框叠成，
矢量图元里没有引线端点信息，无法判定哪条标注归哪条连线；本讲义只按「两两配对」的通则排列，
不指定斜边归属。}

\gapitem{[待确认]（p32）}{两张橙色曲线图（闭合缠绕的环状与「$\infty$」形轨迹）横纵轴均无轴名、无刻度，
全页除页眉外无说明文字，无法确定画的是哪两个变量。本讲义只把它作为「非线性／混沌」的插图引用，
不解释其坐标含义。}

\subhdr{4.3 四步判符号法（可执行）}

\callout{\textbf{见到一个元件，按这四步走，不会错}：\\
\textbf{第 1 步} 在图上找出电压参考极性的「$+$」端（有 $+/-$ 直接读；只有箭头就取箭尾；只有 $U_{AB}$ 就取 $A$）。\\
\textbf{第 2 步} 看电流参考方向的箭头：是\textbf{流入}这个「$+$」端，还是\textbf{流出}它。\\
\textbf{第 3 步} 流入 $\Rightarrow$ 关联 $\Rightarrow$ 用 $u=Ri$、$u=L\,\mathrm{d}i/\mathrm{d}t$、$i=C\,\mathrm{d}u/\mathrm{d}t$。\\
\textbf{第 4 步} 流出 $\Rightarrow$ 非关联 $\Rightarrow$ 三式\textbf{全部}加负号，写成 $u=-Ri$ 等。\\
最后报答案时，负号要么翻译成「实际方向与图上相反」，要么在图上把箭头翻过来再擦掉负号，二选一，不许含糊。}

\subhdr{4.4 为什么这一节是整册的枢纽}

后面三节的全部内容，都是 §4 这条符号约定的直接推论：

\par\vspace{2pt}
\noindent\begin{tabular}{@{}p{3.6cm}p{11.0cm}@{}}
\tabtop
\thead{后面要讲的} & \thead{它就是 §4 的哪一句}
\\
\tabmid
功率的正负（§5） & $p = ui$：关联方向下才算「吸收」，符号即结论（p37 的判据其实只有这一句）
\\
两种理想源（§6） & 电源就是「$u$ 与 $i$ 不独立」的元件：电压源钉住 $u$，电流源钉住 $i$
\\
等效变换（§7） & 等效的定义是「端口的 $u$--$i$ 关系不变」，而 $u$--$i$ 关系的写法就是关联方向下的那一式
\tabbot
\end{tabular}

%=============================================================
\hdr{§5}{功率：$P$ 的正负与「吸收／发出」（p37--p43）}

\subhdr{5.1 p37 的判据，补上它漏掉的三处}

p37 的原文是两句话加一个式子：

\par\medskip\noindent
「Notice: Not all the sources supply power. How to find out if a source provide power or absorb power?
$P = U\cdot I$. If $U$ and $I$ are in the same direction, then it will absorb power.
If $U$ and $I$ are in the opposite direction, then it will supply power.」
\par\medskip

\begin{center}
\begin{tikzpicture}[scale=0.8,>=Stealth]
  \draw[cwire] (0,2) -- (1.6,2);
  \draw[cres] (1.6,1.5) rectangle (4.0,2.5);
  \node[cval] at (2.8,2.0) {source};
  \draw[cwire] (4.0,2) -- (5.6,2);
  \node[cval,red] at (1.75,2.75) {$+$};
  \node[cval,red] at (1.75,1.15) {$-$};
  \node[cval,red] at (1.2,2.5) {$U$};
  \draw[carr,blue] (0.3,2.6) -- (1.45,2.6);
  \node[cval,blue] at (0.85,3.1) {$I$};
\end{tikzpicture}
\end{center}

话本身不错，但有\textbf{三处}必须补，否则一定用错：

\par\vspace{2pt}
\noindent\begin{tabular}{@{}p{0.8cm}p{5.4cm}p{8.6cm}@{}}
\tabtop
\thead{漏处} & \thead{原文的问题} & \thead{补法（本讲义）}
\\
\tabmid
1 & 「same direction」指真实方向还是参考方向？ & 只能是\textbf{参考方向}。
真实方向恰恰是待求的，若要先知道真实方向才能判功率，这条判据就没有用了
\\
2 & 只给「同向吸收、反向发出」的定性表述 & 补上\textbf{代数判据}（下面的定理），
让符号本身成为结论，不必再回图上看方向
\\
3 & 把「叫电源」与「发功率」混在一起 & 器件叫 source 是因为它维持电压或电流，
\textbf{不代表此刻在发出功率}；它在充电时就是负载
\tabbot
\end{tabular}

\subhdr{定理（功率的符号判据）}

设某二端元件的端电压为 $u$、电流为 $i$，且 $u$、$i$ 取\textbf{关联}参考方向，令

\[ p = ui . \]

则 $p$ 是元件\textbf{吸收}的功率：$p>0$ 表示确实吸收，$p<0$ 表示实际在发出。
若 $u$、$i$ 取非关联方向，则同一个 $ui$ 表示\textbf{发出}的功率（等价地，吸收功率为 $-ui$）。

\textbf{为什么}：吸收功率的定义是「电场力对电荷做功、电荷把能量交给元件」。
在关联方向下，正电荷由高电位的 $+$ 端流入、由低电位的 $-$ 端流出，每单位电荷交出的能量正是
$u$，单位时间流过的电荷是 $i$，故单位时间交出的能量为 $ui$。方向非关联时电荷实际是「从低电位端流入」，
能量的流向随之反向，符号也要反。（这就是 §4 的符号规则在「能量」上的同一句话。）

\aside{由此得到最好用的一个自查：对整个电路把每个元件的 $p=ui$（都按关联方向）加起来，
必须\textbf{等于零}，即 $\sum p = 0$。下面 §5.2 的表最后一列就是这个检查，四行全为零。}

\subhdr{5.2 四张 3A/3V 图：p38--p41 逐图算清}

p38--p41 是同一族电路：左边一支 $3\,\mathrm{A}$ 理想电流源（箭头向上），右边一支 $3\,\mathrm{V}$
理想电压源（$+$ 在上），中间并一支电阻 $R$。p38 没有 $R$，p39、p40、p41 依次是
$R=1\,\Omega$、$0.5\,\Omega$、$2\,\Omega$。四页\textbf{都没有给结论}，只反复问
「Find out which one supply power, which one absorb power?」。

\begin{center}
\begin{tikzpicture}[scale=0.78,>=Stealth]
  % 并联组
  \draw[cwire] (1.6,0.2) -- (1.6,1.1);
  \draw[csrc] (1.6,1.55) circle (0.42);
  \draw[thick] (1.18,1.55) -- (2.02,1.55);
  \draw[carr,red] (1.6,2.1) -- (1.6,2.75);
  \draw[cwire] (1.6,1.97) -- (1.6,3.3);
  \node[cval] at (0.85,1.55) {$3\,\mathrm{A}$};
  \draw[cwire] (1.6,3.3) -- (7.2,3.3);
  \draw[cwire] (1.6,0.2) -- (7.2,0.2);
  % 电阻
  \draw[cwire] (4.0,3.3) -- (4.0,2.9);
  \draw[cres] (3.72,1.8) rectangle (4.28,2.9);
  \draw[cwire] (4.0,1.8) -- (4.0,0.2);
  \node[cval] at (4.55,2.35) {$R$};
  \draw[carr,blue] (4.0,2.85) -- (4.0,2.0);
  % 电压源
  \draw[cwire] (7.2,3.3) -- (7.2,2.3);
  \draw[csrc] (7.2,1.75) circle (0.42);
  \draw[cwire] (7.2,1.33) -- (7.2,0.2);
  \node[cval] at (7.75,1.75) {$3\,\mathrm{V}$};
  \node[cval,red] at (7.9,2.8) {$+$};
  \node[cval,red] at (7.9,0.75) {$-$};
  \node[cval] at (4.4,-0.3) {p39--p41：三支并联（p38 无 $R$）};
\end{tikzpicture}
\hspace{1.2cm}
\begin{tikzpicture}[scale=0.78,>=Stealth]
  % 串联回路
  \draw[cwire] (0.6,0.6) -- (0.6,1.3);
  \draw[csrc] (0.6,1.75) circle (0.42);
  \draw[thick] (0.18,1.75) -- (1.02,1.75);
  \draw[carr,red] (0.6,2.3) -- (0.6,2.95);
  \draw[cwire] (0.6,2.17) -- (0.6,3.3);
  \node[cval] at (-0.2,1.75) {$3\,\mathrm{A}$};
  \draw[cwire] (0.6,3.3) -- (5.2,3.3);
  \draw[cres] (2.6,3.02) rectangle (3.4,3.58);
  \node[cval] at (3.0,4.1) {$1\,\Omega$};
  \draw[carr,blue] (1.5,4.0) -- (2.3,4.0);
  \node[cval,blue] at (1.9,4.45) {$i$};
  \draw[cwire] (5.2,3.3) -- (5.2,2.3);
  \draw[csrc] (5.2,1.75) circle (0.42);
  \draw[cwire] (5.2,1.33) -- (5.2,0.6);
  \draw[cwire] (0.6,0.6) -- (5.2,0.6);
  \node[cval] at (5.75,1.75) {$3\,\mathrm{V}$};
  \node[cval,red] at (5.9,2.85) {$+$};
  \node[cval,red] at (5.9,0.7) {$-$};
  \node[cval] at (2.9,-0.05) {p42：单回路串联};
\end{tikzpicture}
\end{center}

先算 p38--p41。把上方导线记作节点 $T$、下方记作 $B$，取 $V_B=0$。
理想电压源钉住 $V_T$，理想电流源钉住注入 $T$ 的电流，于是\textbf{全部未知量只有一个}：
流过电压源支路的电流。用 KCL 于 $T$：

\[ 3 = I_R + I_V, \qquad I_R = \frac{3\,\mathrm{V}}{R}, \qquad
I_V = 3 - \frac{3}{R}\ \ (\text{向下为正，即流入 } T \text{ 端}). \]

\begin{center}
\begin{tabular}{@{}p{1.0cm}p{1.6cm}p{1.9cm}p{2.1cm}p{2.4cm}p{2.3cm}p{1.9cm}p{1.0cm}@{}}
\tabtop
\thead{页} & \thead{$R$} & \thead{$I_R$（向下）} & \thead{$I_V$（向下）}
& \thead{3V 源} & \thead{3A 源} & \thead{$R$ 耗散} & \thead{$\sum p$}
\\
\tabmid
p38 & 无 & $0$ & $3\,\mathrm{A}$ & 吸收 $9\,\mathrm{W}$ & 发出 $9\,\mathrm{W}$ & --- & $0$
\\
p39 & $1\,\Omega$ & $3\,\mathrm{A}$ & $0$ & $0$ & 发出 $9\,\mathrm{W}$ & 吸收 $9\,\mathrm{W}$ & $0$
\\
p40 & $0.5\,\Omega$ & $6\,\mathrm{A}$ & $-3\,\mathrm{A}$ & 发出 $9\,\mathrm{W}$ & 发出 $9\,\mathrm{W}$ & 吸收 $18\,\mathrm{W}$ & $0$
\\
p41 & $2\,\Omega$ & $1.5\,\mathrm{A}$ & $1.5\,\mathrm{A}$ & 吸收 $4.5\,\mathrm{W}$ & 发出 $9\,\mathrm{W}$ & 吸收 $4.5\,\mathrm{W}$ & $0$
\tabbot
\end{tabular}
\end{center}

\aside{表的读法：$I_V>0$ 表示电流向下流过电压源、\textbf{从 $+$ 端进入}，故吸收；$I_V<0$ 表示实际向上、
\textbf{从 $+$ 端流出}，故发出。$3\,\mathrm{A}$ 源那一列恒为「发出 $9\,\mathrm{W}$」，因为它的电流总是
$3\,\mathrm{A}$ 离开 $+$ 端、端电压恒为 $3\,\mathrm{V}$。$R$ 一列按 $P=U^2/R$ 算，与 $I_R^2R$ 同值。}

\subhdr{这张表讲了四件课件的文字讲不清的事}

\par\vspace{2pt}
\noindent\begin{tabular}{@{}p{0.8cm}p{13.6cm}@{}}
\tabtop
\thead{序} & \thead{结论}
\\
\tabmid
1 & \textbf{同一个 $3\,\mathrm{V}$ 源可以从「吸收 9 W」变成「发出 9 W」}，只因为它并联的电阻从无到有、
从 $0.5\,\Omega$ 变到 $2\,\Omega$。可见「某个器件是电源还是负载」不是器件属性，是\textbf{工作点属性}
\\
2 & $R=1\,\Omega$ 时电压源电流恰为零（表格第二行）——这不是巧合：$3\,\mathrm{V}/1\,\Omega = 3\,\mathrm{A}$，
电阻把电流源送来的电流全部接走了。这是「电流源与电压源同值配平」的临界点
\\
3 & p37 那句「Not all the sources supply power」在 p40 得到最强演示：\textbf{两个源同时在发出}，
而 p41 又变成「一个发、一个吸」。光看器件名字永远判不出来
\\
4 & 四行 $\sum p = 0$，能量守恒自动成立。这张表的每一行都是一次完整验算，比背结论可靠
\tabbot
\end{tabular}

\subhdr{5.3 p42--p43：拓扑悄悄换了，答案要重算}

p42 用的是和 p41 完全相同的三个元件值（$3\,\mathrm{A}$、$1\,\Omega$、$3\,\mathrm{V}$），
但 $1\,\Omega$ 从\textbf{竖直支路}挪到了\textbf{水平顶边}，三个元件变成一条\textbf{无分支的串联回路}。
按矢量图元核对：p39 的电阻矩形是竖的（宽 $20$、高 $64$ pt），p42 的电阻矩形是横的（宽 $61$、高 $20$ pt）。
课件没有一句话说明这次切换，这是本册第三处「拼接缺口」。

串联情形下，$3\,\mathrm{A}$ 电流源仍强制循环电流 $3\,\mathrm{A}$：由左支路向上、沿顶边向右过 $1\,\Omega$、
再沿右支路向下过 $3\,\mathrm{V}$ 源。沿电流方向绕行一圈，把各段的\textbf{电位变化}相加为零
（升高记正、降低记负）：

\[ u_{cs} - 3\times1 - 3 = 0 \quad\Longrightarrow\quad u_{cs} = 6\,\mathrm{V}, \]

其中 $u_{cs}$ 是电流源上端相对下端的电压（参考极性取「$+$ 在上」，与它的箭头一致）。
于是三个元件的功率是

\par\vspace{1pt}
\noindent
电流源：电流离开 $+$ 端，故 $\ p = 6\times(-3) = -18\,\mathrm{W}$（吸收），即\textbf{发出} $18\,\mathrm{W}$；\\
$1\,\Omega$ 电阻：$3^2\times1 = 9\,\mathrm{W}$ \textbf{吸收}；\\
$3\,\mathrm{V}$ 电压源：电流从 $+$ 端流入，故\textbf{吸收} $3\times3 = 9\,\mathrm{W}$。
\par\vspace{1pt}

校验：$18 = 9+9$ $\checkmark$。请注意与 p39 的差别：\textbf{p39 里那个 $1\,\Omega$ 并联不动电压源的电流，
p42 里它串在回路中，反而抬高了电流源的端电压}（从 $3\,\mathrm{V}$ 抬到 $6\,\mathrm{V}$）。

\subhdr{5.3.1 p43 的提问：电流源总是发出功率吗}

p43 在 p42 的图下方加了一句「Does the current source in the circuit always provide power?」，
没有给答案。正确的回答是\textbf{不是}，理由只用一行：

\callout{\textbf{答}：在 p38--p43 这五张图里，这个 $3\,\mathrm{A}$ 源\textbf{都}在发出功率，
但这是\textbf{这五个电路}的性质，不是电流源的普遍性质。电流源的电流由自己固定，
它两端的电压极性却由\textbf{外电路}决定；一旦端电压的参考极性反过来，同一个电流源就吸收功率。\\
\textbf{验证}：把 p39 的 $3\,\mathrm{V}$ 源反接（$+$ 端改在下）。顶节点变成 $-3\,\mathrm{V}$。
$3\,\mathrm{A}$ 源的电流仍离开上端（它的 $+$ 端），而 $u$ 由 $+3\,\mathrm{V}$ 变成 $-3\,\mathrm{V}$，
故 $p$ 由 $-9\,\mathrm{W}$ 变成 $+9\,\mathrm{W}$：它\textbf{吸收} $9\,\mathrm{W}$。
同时 $1\,\Omega$ 上的电流反向，仍吸收 $9\,\mathrm{W}$；$3\,\mathrm{V}$ 源改为发出 $18\,\mathrm{W}$。
三件相加 $9+9-18=0$ $\checkmark$。}

\warnbox{\textbf{一个必须掰清楚的措辞}：p37 的结论「If $U$ and $I$ are in the opposite direction,
then it will supply power」里的「it」是\textbf{这个元件}，不是「电源这一类器件」。
电器件名与功率方向之间没有必然联系：发电机在正常发电时发出功率，在电动机拖动的「发电机工况」下吸收功率；
电池放电时发出、充电时吸收。判据永远只是 §5.1 的 $p=ui$ 与它的符号。}

%=============================================================
\hdr{§6}{电源模型：两种理想源与两种实际源（p33--p36、p44--p45）}

\subhdr{6.1 理想电压源：钉住电压，电流随你（p33--p34）}

p33 的特性框只有两条：

\par\medskip\noindent
（1）The output voltage is constant and its value is constant equal to electromotive force.\\
（2）The current in the voltage source is determined by the external circuit.
\par\medskip

第一句说的是「$U=U_S$ 与电流无关」，第二句说的是「电流由外电路决定」——两句合起来才是完整定义：
\textbf{理想电压源是一个把端电压钉死、而对流过它的电流不设限的元件}。
p33 右侧的伏安特性是一条水平的直线 $U=U_S$，p33 左侧的符号图标着 $U=\text{const.}$，都是这句话的图形版。

它为什么会这样？因为它\textbf{没有内阻}。§1.6 的表里 $U=U_S\dfrac{R}{R_0+R}$，
令 $R_0\to0$ 就得到 $U=U_S$，与 $R$ 无关；而 $I=\dfrac{U_S}{R_0+R}\to\dfrac{U_S}{R}$，完全由负载决定。

p34 用一组数把第二条钉实：$U_S=10\,\mathrm{V}$、$R_1=R_2=2\,\Omega$。

\par\vspace{2pt}
\noindent\begin{tabular}{@{}p{4.2cm}p{4.0cm}p{7.4cm}@{}}
\tabtop
\thead{外电路} & \thead{电流} & \thead{说明}
\\
\tabmid
只接 $R_1$ & $I=10/2=5\,\mathrm{A}$ & 负载 $2\,\Omega$
\\
$R_1\parallel R_2$ & $I=10/1=10\,\mathrm{A}$ & 负载变成 $1\,\Omega$，电流翻倍，
而端电压仍是 $10\,\mathrm{V}$
\tabbot
\end{tabular}

\warnbox{\textbf{两条禁令（课件没写，但 p38 立刻就用到了它的变体）}：\\
（1）\textbf{理想电压源不能短路}：$R=0$ 时 $I=U_S/R_0\to\infty$，物理上做不到。
p16 的「短路电流最大」说的正是这个极限，真实电路靠 $R_0$ 兜住。\\
（2）\textbf{两个电压值不同的理想电压源不能并联}：它们的端电压被各自钉死，
并联后同一个节点要同时等于两个不同的数，方程无解。（p38 那种「$3\,\mathrm{A}$ 源与 $3\,\mathrm{V}$ 源并联」
之所以有解，是因为电流源不钉电压，只钉电流。）}

\subhdr{6.1.1 一条极有用的推论：与理想电压源并联的元件「不影响其余电路」}

\callout{\textbf{推论（本讲义补，§8 的 E1 要用）}：若某元件与一个\textbf{理想}电压源并联，
则去掉这个元件，其余各点的电压与电流\textbf{全都不变}；变的只是「这个理想源自己要提供多少电流」。\\
\textbf{为什么}：并联意味着两者端电压相同，而端电压已被理想源钉死，故该元件的存在与否
不改变任何节点电压；节点电压不变 $\Rightarrow$ 其余支路电流不变。\\
\textbf{注意}：这条只在\textbf{理想}源成立。若源有内阻，那个并联元件会分走电流、改变内阻上的压降，
端电压就变了。}

\subhdr{6.2 理想电流源：钉住电流，电压随你（p35--p36）}

p35 的两条特性与 p33 恰好对偶：

\par\medskip\noindent
（1）The output current is constant and its value is constant equal to the current source current.\\
（2）The output voltage is determined by the external circuit.
\par\medskip

p35 的伏安特性画在 $U$--$I$ 平面上（横轴 $U$、纵轴 $I$）是一条水平线 $I=I_S$；
p36 用一组数演示第二条：$I_S=1\,\mathrm{A}$，$R=1\,\Omega$ 时 $U=1\,\mathrm{V}$，
$R=10\,\Omega$ 时 $U=10\,\mathrm{V}$——电流不变，电压跟着外电路走。

对偶地，也有两条禁令与一条推论：

\par\vspace{2pt}
\noindent\begin{tabular}{@{}p{0.8cm}p{13.6cm}@{}}
\tabtop
\thead{序} & \thead{内容}
\\
\tabmid
禁令 1 & \textbf{理想电流源不能开路}：开路时 $R\to\infty$，$U=I_SR\to\infty$
\\
禁令 2 & \textbf{两个电流值不同的理想电流源不能串联}：同一个支路要被两个源钉成两个不同的电流
\\
推论 & \textbf{与理想电流源串联的元件「不影响其余电路」}：串联意味着两者电流相同，
电流已被钉死，故该元件只分走电压，不改变注入节点的电流（§8 的 E1 要用）
\tabbot
\end{tabular}

\subhdr{6.3 实际电压源：课件只画了曲线，没给方程（p44）}

p44 把 Ideal 与 Practical 两栏并排：左栏电压源端电压恒定、$U$--$I$ 是一条水平线；
右栏在电压源符号下面多串了一个电阻（这就是内阻 $R_0$），$U$--$I$ 变成一条\textbf{向右下方倾斜}的直线。
课件\textbf{没有写这条直线是哪个方程}，只写了「Practical／$U\ne$ const.」。补上：

\subhdr{实际电压源的端口方程（本讲义补）}

设理想电动势为 $U_S$、内阻为 $R_0$，电流 $I$ 由 $+$ 端流出（对外供电）。则端电压为

\[ U = U_S - I R_0 . \]

\textbf{推导}：内阻 $R_0$ 与理想电动势 $U_S$ 串联。对外供电时电流 $I$ 从 $+$ 端流出、经外电路回到 $-$ 端，
故在电源\textbf{内部}它是由 $-$ 端流向 $+$ 端。从 $-$ 端沿电源内部走到 $+$ 端：
经过理想电动势时电位\textbf{升高} $U_S$；经过内阻时，因为电流的流向与行走方向相同，
电位\textbf{下降} $IR_0$。于是端电压

\[ U = V_+ - V_- = \underbrace{U_S}_{\text{电动势升高}} - \underbrace{IR_0}_{\text{内阻压降}} . \]

它就是一条直线：纵截距 $U_S$（空载电压），横截距 $U_S/R_0$（短路电流），斜率 $-R_0$
（所以内阻越大，曲线越陡）。三个特殊点与 §1.6 的表完全对应：

\par\vspace{2pt}
\noindent\begin{tabular}{@{}p{3.0cm}p{3.0cm}p{3.6cm}p{5.4cm}@{}}
\tabtop
\thead{工作点} & \thead{$I$} & \thead{$U$} & \thead{在 p44 右栏曲线上的位置}
\\
\tabmid
开路 & $0$ & $U_S$ & 纵轴上的截距
\\
工作点 & $U_S/(R_0+R)$ & $U_S-IR_0$ & 截距之间的某一点
\\
短路 & $U_S/R_0$ & $0$ & 横轴上的截距
\tabbot
\end{tabular}

\begin{center}
\begin{tikzpicture}[scale=0.9,>=Stealth]
  % 理想
  \draw[->,>=Stealth,thick] (0,0) -- (3.6,0) node[right,clab] {$I$};
  \draw[->,>=Stealth,thick] (0,0) -- (0,2.6) node[above,clab] {$U$};
  \draw[very thick,blue] (0,2.0) -- (3.3,2.0);
  \node[clab,left] at (0,2.0) {$U_S$};
  \node[clab] at (1.9,-0.55) {理想：$U=U_S$};
  % 实际
  \draw[->,>=Stealth,thick] (5.6,0) -- (9.2,0) node[right,clab] {$I$};
  \draw[->,>=Stealth,thick] (5.6,0) -- (5.6,2.6) node[above,clab] {$U$};
  \draw[very thick,red!70!black] (5.6,2.0) -- (8.6,0.5);
  \node[clab,left] at (5.6,2.0) {$U_S$};
  \node[clab,below] at (8.6,0.4) {$\frac{U_S}{R_0}$};
  \draw[thick,dashed] (6.6,2.0) -- (6.6,1.5) -- (5.6,1.5);
  \node[clab,left] at (5.6,1.5) {$U$};
  \node[clab] at (7.6,-0.55) {实际：$U=U_S-IR_0$};
\end{tikzpicture}
\end{center}

\subhdr{6.4 实际电流源：对偶的一行（p45）}

p45 的结构与 p44 完全对称：左栏理想电流源的 $I$--$U$ 特性是水平线 $I=I_S$；
右栏在电流源旁边\textbf{并}了一个内阻，输出电流随端电压增大而下降。它的端口方程同样没写，补上：

\subhdr{实际电流源的端口方程（本讲义补）}

设理想电流源的电流为 $I_S$、内阻为 $R_0$（与源\textbf{并联}），则对外输出的电流为

\[ I = I_S - \frac{U}{R_0} . \]

两个截距：空载（$I=0$）时 $U_{\mathrm{OC}} = I_SR_0$；短路（$U=0$）时 $I_{\mathrm{SC}} = I_S$。

\par\vspace{2pt}
\noindent\begin{tabular}{@{}p{2.4cm}p{4.6cm}p{4.6cm}@{}}
\tabtop
\thead{对照项} & \thead{实际电压源（p44）} & \thead{实际电流源（p45）}
\\
\tabmid
内阻接法 & 串 联 & 并 联
\\
端口方程 & $U = U_S-IR_0$ & $I = I_S-U/R_0$
\\
特性曲线 & 下垂直线，纵截距 $U_S$ & 下垂直线，纵截距 $I_S$
\\
空载 & $I=0$，$U=U_S$ & $I=0$，$U=I_SR_0$
\\
短路 & $U=0$，$I=U_S/R_0$ & $U=0$，$I=I_S$
\\
退化为理想源 & $R_0\to0$ & $R_0\to\infty$
\tabbot
\end{tabular}

\warnbox{\textbf{为什么「实际」二字如此重要}：一个\textbf{理想}电压源（$R_0=0$）与一个\textbf{理想}电流源
（$R_0=\infty$）都无法用对方等效——前者的端口方程与 $I$ 无关，后者的与 $U$ 无关，
两条曲线一个是水平线、一个是竖直线，永远对不上。p47 那句「For practical not for ideal」
说的就是这件事：\textbf{等效变换只在 $0<R_0<\infty$ 时才有意义}。}

%=============================================================
\hdr{§7}{电源等效变换：五问、公式与「为什么成立」（p46--p52）}

\subhdr{7.1 五问：课件答了四问，第五问只有演示（p46--p47）}

p46、p47、p48--p52 的左栏始终挂着同一组五个问题，高亮逐个往后移。这种排版是在告诉你：
\textbf{这五个问题就是这一节的骨架}。逐条对照课件给了什么、缺什么：

\par\vspace{2pt}
\noindent\begin{tabular}{@{}p{1.2cm}p{3.9cm}p{5.8cm}p{1.1cm}p{3.4cm}@{}}
\tabtop
\thead{序} & \thead{问题（p46 原文）} & \thead{课件给的回答} & \thead{页} & \thead{缺什么}
\\
\tabmid
（1） & What is ST? & 「Simply means replacing one source by an equivalent source.」
即用一个等效源替换另一个源 & p46 & ---
\\
（2） & What is the Purpose of ST? & 「The purpose is to simplify the calculation.」 & p46 & ---
\\
（3） & Under what conditions can ST be carried out? & 「（1）For practical not for ideal.
（2）For outside not for inside.」（课件这两条用的是带圈数字） & p47 & 两条都要展开成判据（§7.3、§7.4）
\\
（4） & How to transform? & $R'=R$，$I_S=U_S/R$；反向 $U_S=I_SR$ & p48--p50 & 方向怎么对（§7.2）
\\
（5） & Why can such a transformation be made? & \textbf{没有正面回答}，只用 p51--p52
的「$2\,\mathrm{A}$」演示了一次 & p51--p52 & 证明（§7.3）
\tabbot
\end{tabular}

\gapitem{表述缺口}{p47 的第三问答案里有一句「It can not only be used to calculate the parameters of the
source itself.」按字面读是「\textbf{不仅}能算电源自身的参数」，而它上一条刚说「For outside not for inside」，
意思正好相反。从上下文（配合 p54「No influence on external circuits」）判断，
这里想说的是「\textbf{只能}用来算外电路的量，不能用来算电源内部」。课件这句英文的 not only 用坏了，
本讲义按「只能算外部」理解，并在 §7.4 给出它为什么如此。}

\subhdr{7.2 怎么换：公式与「箭头对正端」（p48--p50）}

p48 一张图把两个模型上下对照，给了两条公式（黄色框）：

\[ R' = R, \qquad I_S = \frac{U_S}{R}, \qquad\text{反向：}\quad U_S = I_S R' . \]

p49 给出待换的数值（$10\,\mathrm{V}$ 串 $2\,\Omega$），p50 给出结果（$5\,\mathrm{A}$ 并 $2\,\Omega$），
数字对得上：$I_S = 10/2 = 5\,\mathrm{A}$，$R'=2\,\Omega$。

课件\textbf{只用箭头画了方向}，没有写那条最容易错的规则，补上：

\callout{\textbf{方向规则（本讲义补）}：等效前后，\textbf{电流源的箭头尖端对应电压源的正端}。\\
理由：实际电压源对外供电时，电流从 $+$ 端流出；实际电流源的箭头画的就是「源自己送出电流的方向」，
故箭头也指向 $+$ 端。p48--p50 的图上，$10\,\mathrm{V}$ 源的 $+$ 在上、$5\,\mathrm{A}$ 源的箭头向上，
正是这条规则。\textbf{把箭头画反，短路电流就反号}，是这一节最高频的错误。}

\subhdr{7.3 为什么成立：第五问的正面回答（本讲义补）}

课件用 p51--p52 的打法：给两个模型各接同一个 $3\,\Omega$ 负载，算出两边都是 $2\,\mathrm{A}$，
然后说「看，一样」。这在教学上有效，但\textbf{不是证明}——它只验了一个负载点。
真正的证明只要三行，核心是「等效」的定义：

\callout{\textbf{定义（端口等效）}：两个二端网络，如果对\textbf{任意}外部电路都给出相同的
端电压—端电流关系（相同的端口方程），则称二者等效。\\
注意：等效是对「端口」说的，不是对「内部」说的。这条定义后面决定了一切。}

\subhdr{定理（电源等效变换成立）}

设实际电压源为 $U_S$ 串 $R$，实际电流源为 $I_S$ 并 $R'$。若

\[ R' = R, \qquad I_S = \frac{U_S}{R}, \]

则二者端口等效。

\textbf{证明}：两个模型的端口方程（均取「电流 $I$ 由 $+$ 端流出」为正向）分别是

\[ \text{电压源侧：}\quad U = U_S - IR ; \qquad
   \text{电流源侧：}\quad I = I_S - \frac{U}{R'} . \]

把 $I_S = U_S/R$、$R'=R$ 代入第二式的右边：

\[ I = \frac{U_S}{R} - \frac{U}{R}
  \quad\Longrightarrow\quad
  IR = U_S - U
  \quad\Longrightarrow\quad
  U = U_S - IR , \]

与第一式\textbf{逐字相同}。两个模型给出的端口方程是同一个方程，
故对任意外电路（任意负载电流或电压）都给出同一个端口解，按定义即等效。$\square$

\aside{这个证明的每一步都可以回指课件：$U=U_S-IR$ 是 §6.3 由 p44 的曲线补出的方程，
$I=I_S-U/R'$ 是 §6.4 由 p45 的曲线补出的方程，$R'=R$、$I_S=U_S/R$ 是 p48 黄色框里的原式。
课件缺的只是「把两式并排放、发现它们是同一个」这一步。}

\subhdr{7.3.1 一个更直观的等价证法：两点定一条直线}

两个端口方程都是\textbf{一次}的（直线），所以只要证明两条直线过同样的\textbf{两个点}即可。
取最有辨识度的两点——空载与短路：

\par\vspace{2pt}
\noindent\begin{tabular}{@{}p{2.4cm}p{5.0cm}p{5.0cm}@{}}
\tabtop
\thead{工作点} & \thead{电压源侧} & \thead{电流源侧}
\\
\tabmid
空载（$I=0$） & $U_{\mathrm{OC}} = U_S$ & $U_{\mathrm{OC}} = I_SR' = \dfrac{U_S}{R}\cdot R = U_S$ $\checkmark$
\\
短路（$U=0$） & $I_{\mathrm{SC}} = \dfrac{U_S}{R}$ & $I_{\mathrm{SC}} = I_S = \dfrac{U_S}{R}$ $\checkmark$
\tabbot
\end{tabular}

两点相同、都是直线 $\Rightarrow$ 整条直线相同。这个证法还有个附带好处：
它把 p52 的「$2\,\mathrm{A}$」从「偶然对上」变成了「必然对上」的\textbf{第三个抽查点}：用分流公式独立算一遍，
$I_{3\Omega} = 5\times\dfrac{2}{2+3} = 2\,\mathrm{A}$，与电压源侧的 $10/(2+3) = 2\,\mathrm{A}$ 一致。

\subhdr{7.4 「只对外不对内」到底差多少：一个可算的数字}

p47 的第二条条件「For outside not for inside」在课件里只有一句话，没有数字。
但它是可以算的，而且算出来差得很远。取 p49/p50 那一对模型（$10\,\mathrm{V}$ 串 $2\,\Omega$
$\leftrightarrow$ $5\,\mathrm{A}$ 并 $2\,\Omega$），让它们都处于\textbf{空载}：

\par\vspace{2pt}
\noindent\begin{tabular}{@{}p{3.2cm}p{4.4cm}p{4.4cm}@{}}
\tabtop
\thead{量} & \thead{电压源侧} & \thead{电流源侧}
\\
\tabmid
端口电压 & $10\,\mathrm{V}$ & $5\times2 = 10\,\mathrm{V}$ $\checkmark$ 相同
\\
端口电流 & $0$ & $0$ $\checkmark$ 相同
\\
内阻上的电流 & $0$ & $5\,\mathrm{A}$（全部在电源内部环流）
\\
电源内部损耗 & $0\,\mathrm{W}$ & $5^2\times2 = 50\,\mathrm{W}$ \textbf{不同}
\\
\tabbot
\end{tabular}

外面一毫不差，里面差了 $50\,\mathrm{W}$。这就是「等效」二字的边界：
\textbf{它保的是端口方程，不保内部的电压、电流与功率}。所以
（1）可以用变换去算外电路的电压电流；（2）不可以用变换后的电路去算原来那个源内部的损耗、
也不可以问「变换前那个 $2\,\Omega$ 里的电流是多少」。p47 那句被写坏的英文，正确内容就是这两句。

\warnbox{\textbf{再补一条 p47 没写的适用条件}：变换只对\textbf{电源与电阻的串联／并联组合}成立，
而且变换前后\textbf{不能有第三个节点被合并或分裂}。具体地说，
「$U_S$ 串 $R$」这一整块的两端必须恰好是某两个节点，中间不允许再引出支路；
若中间有引出（例如 p53 的电路），就要先把中间引出处理掉再换——这正是 E1 第一步要做的事（§8.1）。}

%=============================================================
\hdr{§8}{主例题全程：E1（p53--p58）}

\subhdr{8.1 原图（p53）与两处「没有影响」的部分}

p53 标着「E1」，是一张只有元件、没有题面的图；p54 起了两行字：
「（1）The part connected in parallel with the ideal voltage source；
（2）The part connected in series with the ideal current source」，下面一行
「No influence on external circuits」。三句话其实就是解题的第一步：
\textbf{先去掉两个不起作用的元件}。

按该页矢量图元复原，原电路的节点只有四个：上方导线被两个串联电阻分成三段，
分别记作 $T_1$、$T_2$、$T_3$，下方导线全通，记作 $B$。

\begin{center}
\begin{tikzpicture}[scale=0.88,>=Stealth]
  % 底边
  \draw[cwire] (0.6,0) -- (8.2,0);
  % 顶边（三段＋两个串联电阻）
  \draw[cwire] (0.6,5) -- (2.6,5);
  \draw[cres] (2.6,4.72) rectangle (3.4,5.28);
  \node[cval] at (3.0,5.8) {$2\,\mathrm{k}\Omega$};
  \draw[cwire] (3.4,5) -- (6.6,5);
  \draw[cres] (6.6,4.72) rectangle (7.4,5.28);
  \node[cval] at (7.0,5.8) {$1\,\mathrm{k}\Omega$};
  \draw[cwire] (7.4,5) -- (8.2,5);
  % 5V 源（x=0.6）
  \draw[cwire] (0.6,0) -- (0.6,2.15);
  \draw[csrc] (0.6,2.5) circle (0.35);
  \draw[cwire] (0.6,2.85) -- (0.6,5);
  \node[cval,red] at (1.0,3.1) {$+$};
  \node[cval,red] at (1.0,1.9) {$-$};
  \node[cval,left] at (0.2,2.5) {$5\,\mathrm{V}$};
  % 2kΩ（x=2.0，与 5V 源并联）
  \draw[cwire] (2.0,0) -- (2.0,1.9);
  \draw[cres] (1.72,1.9) rectangle (2.28,3.1);
  \draw[cwire] (2.0,3.1) -- (2.0,5);
  \node[cval,right] at (2.35,2.5) {$2\,\mathrm{k}\Omega$};
  % 1kΩ（x=4.4）
  \draw[cwire] (4.4,0) -- (4.4,1.9);
  \draw[cres] (4.12,1.9) rectangle (4.68,3.1);
  \draw[cwire] (4.4,3.1) -- (4.4,5);
  \node[cval,right] at (4.75,2.5) {$1\,\mathrm{k}\Omega$};
  % 2mA 源 + 串联 2kΩ（x=5.8）
  \draw[cwire] (5.8,0) -- (5.8,1.3);
  \draw[cres] (5.52,1.3) rectangle (6.08,2.5);
  \node[cval,right] at (6.15,1.9) {$2\,\mathrm{k}\Omega$};
  \draw[cwire] (5.8,2.5) -- (5.8,3.4);
  \draw[csrc] (5.8,3.75) circle (0.35);
  \draw[thick] (5.46,3.75) -- (6.14,3.75);
  \draw[carr,red] (5.8,4.15) -- (5.8,4.62);
  \draw[cwire] (5.8,4.1) -- (5.8,5);
  \node[cval,left] at (5.35,3.75) {$2\,\mathrm{mA}$};
  % 2kΩ＋I（x=8.2）
  \draw[cwire] (8.2,0) -- (8.2,1.9);
  \draw[cres] (7.92,1.9) rectangle (8.48,3.1);
  \draw[cwire] (8.2,3.1) -- (8.2,5);
  \node[cval,right] at (8.55,2.5) {$2\,\mathrm{k}\Omega$};
  \draw[carr,blue] (8.2,4.55) -- (8.2,3.55);
  \node[cval,blue] at (8.7,4.15) {$I$};
  % 节点标注
  \fill (0.6,5) circle (1.4pt); \fill (2.0,5) circle (1.4pt);
  \fill (4.4,5) circle (1.4pt); \fill (5.8,5) circle (1.4pt); \fill (8.2,5) circle (1.4pt);
  \foreach \x in {0.6,2.0,4.4,5.8,8.2} { \fill (\x,0) circle (1.4pt); }
  \node[clab] at (1.3,5.4) {$T_1$};
  \node[clab] at (4.9,5.4) {$T_2$};
  \node[clab] at (7.85,5.4) {$T_3$};
  \node[clab] at (4.6,-0.45) {$B$（全底边同一节点）};
\end{tikzpicture}
\end{center}

\par\vspace{2pt}
\noindent\begin{tabular}{@{}p{2.9cm}p{3.0cm}p{6.0cm}p{3.0cm}@{}}
\tabtop
\thead{支路} & \thead{两端} & \thead{元件} & \thead{p54 的处理}
\\
\tabmid
左竖支路 & $T_1$--$B$ & $5\,\mathrm{V}$ 理想电压源，$+$ 在 $T_1$ & 保留（要换）
\\
第二竖支路 & $T_1$--$B$ & $2\,\mathrm{k}\Omega$ & \textbf{去掉}（与理想源并联）
\\
顶横左段 & $T_1$--$T_2$ & $2\,\mathrm{k}\Omega$ & 保留（是 $5\,\mathrm{V}$ 源的串联电阻）
\\
第三竖支路 & $T_2$--$B$ & $1\,\mathrm{k}\Omega$ & 保留
\\
第四竖支路 & $T_2$--$B$ & $2\,\mathrm{mA}$ 电流源（箭头向上）串 $2\,\mathrm{k}\Omega$ & 去掉$2\,\mathrm{k}\Omega$（与理想源串联）
\\
顶横右段 & $T_2$--$T_3$ & $1\,\mathrm{k}\Omega$ & 保留
\\
右竖支路 & $T_3$--$B$ & $2\,\mathrm{k}\Omega$，电流 $I$ 向下 & 保留（待求量所在支路）
\tabbot
\end{tabular}

\begin{center}
\begin{tikzpicture}[scale=0.88,>=Stealth]
  % 简化后的电路
  \draw[cwire] (0.4,0) -- (5.9,0);
  \draw[cwire] (0.4,5) -- (2.4,5);
  \draw[cres] (2.4,4.72) rectangle (3.2,5.28);
  \node[cval] at (2.8,5.8) {$1\,\mathrm{k}\Omega$（原顶横右段）};
  \draw[cwire] (3.2,5) -- (5.9,5);
  % 4.5mA 源
  \draw[cwire] (0.4,0) -- (0.4,1.6);
  \draw[csrc] (0.4,1.95) circle (0.35);
  \draw[thick] (0.06,1.95) -- (0.74,1.95);
  \draw[carr,red] (0.4,2.35) -- (0.4,2.85);
  \draw[cwire] (0.4,2.3) -- (0.4,5);
  \node[cval,left] at (-0.05,1.95) {$4.5\,\mathrm{mA}$};
  % 2/3 kΩ
  \draw[cwire] (1.5,0) -- (1.5,1.9);
  \draw[cres] (1.22,1.9) rectangle (1.78,3.1);
  \draw[cwire] (1.5,3.1) -- (1.5,5);
  \node[cval,right] at (1.85,2.5) {$\frac{2}{3}\,\mathrm{k}\Omega$};
  % 2kΩ 与 I
  \draw[cwire] (5.9,0) -- (5.9,1.9);
  \draw[cres] (5.62,1.9) rectangle (6.18,3.1);
  \draw[cwire] (5.9,3.1) -- (5.9,5);
  \node[cval,right] at (6.25,2.5) {$2\,\mathrm{k}\Omega$};
  \draw[carr,blue] (5.9,4.55) -- (5.9,3.55);
  \node[cval,blue] at (6.4,4.15) {$I$};
  \fill (0.4,5) circle (1.4pt); \fill (1.5,5) circle (1.4pt); \fill (5.9,5) circle (1.4pt);
  \fill (0.4,0) circle (1.4pt); \fill (1.5,0) circle (1.4pt); \fill (5.9,0) circle (1.4pt);
\end{tikzpicture}
\end{center}

\subhdr{8.1.1 这两处为什么可以去掉（补 p54 的证明）}

p54 只给了结论「No influence on external circuits」，没给理由。理由就是 §6.1.1 与 §6.2 的两条推论：

\par\vspace{2pt}
\noindent\begin{tabular}{@{}p{0.8cm}p{12.9cm}@{}}
\tabtop
\thead{序} & \thead{为什么去掉不改变外电路}
\\
\tabmid
（1） & 第二竖支路的 $2\,\mathrm{k}\Omega$ 与\textbf{理想}电压源并联。并联意味着两者端电压相同，
而端电压被理想源钉死为 $5\,\mathrm{V}$。故这个电阻取多少电流都不改变任何节点电压；
它唯一改变的是「$5\,\mathrm{V}$ 源自己要送出多少电流」。其余各点电压电流一律不变。
\\
（2） & 第四竖支路的 $2\,\mathrm{k}\Omega$ 与\textbf{理想}电流源串联。串联意味着两者电流相同，
而电流被理想源钉死为 $2\,\mathrm{mA}$。故这个电阻只分走一段电压，不改变注入 $T_2$ 的电流。
\\
\tabbot
\end{tabular}

\warnbox{\textbf{两条推论的分量在「理想」二字上}：若那个 $5\,\mathrm{V}$ 源带内阻（即实际源），
与它并联的 $2\,\mathrm{k}\Omega$ 会分走电流、改变内阻上的压降，端电压就变了，不能去掉。
同理，与电流源串联的若是\textbf{另一个电流源}，两个源串联要求电流相等，一旦数值不同就无解，
也不是「去掉就好」。p54 的两条规则是理想源的专利。}

\subhdr{8.2 第一步：$5\,\mathrm{V}$ 串 $2\,\mathrm{k}\Omega$ 换成 $2.5\,\mathrm{mA}$ 并 $2\,\mathrm{k}\Omega$（p55--p56）}

去掉两处之后，$T_1$ 只剩下「$5\,\mathrm{V}$ 源串 $2\,\mathrm{k}\Omega$ 到 $T_2$」这一条路，
正好是 §7 的标准形态。按 $I_S=U_S/R$：

\[ I_S = \frac{5\,\mathrm{V}}{2\,\mathrm{k}\Omega} = 2.5\,\mathrm{mA}, \qquad R' = 2\,\mathrm{k}\Omega . \]

于是 $T_1$ 这个节点消失了，$T_2$ 与 $B$ 之间多出两支并联：$2.5\,\mathrm{mA}$ 电流源（箭头向上，
即注入 $T_2$）与 $2\,\mathrm{k}\Omega$ 电阻。p56 的图上这两支被淡黄框圈在一起，旁边写着
「Move them together」——圈的就是这对「源 $+$ 内阻」。

\aside{方向核对：原来 $5\,\mathrm{V}$ 源的 $+$ 端在 $T_1$，电流从 $+$ 端流出、经 $2\,\mathrm{k}\Omega$ 流向 $T_2$，
故等效电流源的箭头应指向 $T_2$（向上）。p56 的 $5/2\,\mathrm{mA}$ 箭头正是向上的，
与本讲义 §7.2 的「箭头尖端对正端」一致。}

\subhdr{8.3 第二步：两个电流源并联合并、两个电阻并联合并（p57）}

此时 $T_2$ 与 $B$ 之间挂着四支并联：$2.5\,\mathrm{mA}$、$2\,\mathrm{k}\Omega$、$1\,\mathrm{k}\Omega$、$2\,\mathrm{mA}$。
同向并联的电流源可以直接相加，并联电阻按并联公式合并：

\[ I_{\text{合}} = 2.5+2 = 4.5\,\mathrm{mA}, \qquad
   R_{\text{合}} = 2\,\mathrm{k}\Omega \parallel 1\,\mathrm{k}\Omega
   = \frac{2\times1}{2+1} = \frac{2}{3}\,\mathrm{k}\Omega . \]

p57 的图上两支电流源（$2.5\,\mathrm{mA}$、$2\,\mathrm{mA}$）与两支电阻（$2\,\mathrm{k}\Omega$、$1\,\mathrm{k}\Omega$）
并排摆开，底下写「When the two current source connected in parallel」，
就是在等这一步相加。\textbf{前提是两支源的方向一致}（都注入 $T_2$），
若方向相反，相加要带负号。

\subhdr{8.4 第三步：分流求 $I$（p58）}

化简到最后，电路是：$4.5\,\mathrm{mA}$ 源与 $\frac{2}{3}\,\mathrm{k}\Omega$ 并联，接在 $T_2$ 与 $B$ 之间；
再从 $T_2$ 串一个 $1\,\mathrm{k}\Omega$ 到 $T_3$，$T_3$ 经 $2\,\mathrm{k}\Omega$ 回到 $B$，$I$ 就流在这条支路上。
于是 $4.5\,\mathrm{mA}$ 要在「$\frac{2}{3}\,\mathrm{k}\Omega$」与「$1+2=3\,\mathrm{k}\Omega$」两条支路之间分流。
p58 左下角给了分流公式，套进去：

\[ I = 4.5\,\mathrm{mA}\times
   \frac{\frac{2}{3}}{\frac{2}{3}+3}
   = 4.5\times\frac{\frac{2}{3}}{\frac{11}{3}}
   = 4.5\times\frac{2}{11}
   = \frac{9}{11}\,\mathrm{mA} \approx 0.818\,\mathrm{mA} . \]

\textbf{课件到这里就停了}：p58 只给出分流公式，\textbf{没有写出 $I$ 的数值}，p59 直接跳到复习题。
这是本册最后一处「答案形态」缺口，上面的 $\frac{9}{11}\,\mathrm{mA}$ 是本讲义补的。

\subhdr{8.4.1 独立复核：直接对原电路做节点分析}

化简对不对，最好用另一种方法验一遍。对 p53 的原电路取 $V_B=0$、$V_{T_1}=5\,\mathrm{V}$，
对 $T_2$、$T_3$ 写 KCL（单位 V、mA、k$\Omega$）：

\[
\begin{cases}
\dfrac{5-V_{T_2}}{2} + 2 = \dfrac{V_{T_2}}{1} + \dfrac{V_{T_2}-V_{T_3}}{1} \\[8pt]
\dfrac{V_{T_2}-V_{T_3}}{1} = \dfrac{V_{T_3}}{2}
\end{cases}
\]

第二式给出 $V_{T_3} = \frac{2}{3}V_{T_2}$；代入第一式并整理得

\[ \frac{5}{2}+2 = \left(1+1-\frac{2}{3}\right)V_{T_2}+\frac{V_{T_2}}{2}
   \;\Longrightarrow\; 4.5 = \frac{11}{6}V_{T_2}
   \;\Longrightarrow\; V_{T_2} = \frac{27}{11}\,\mathrm{V} \approx 2.45\,\mathrm{V}, \]

\[ V_{T_3} = \frac{2}{3}\times\frac{27}{11} = \frac{18}{11}\,\mathrm{V}, \qquad
   I = \frac{V_{T_3}}{2} = \frac{9}{11}\,\mathrm{mA} . \]

与化简法的结果\textbf{逐位相同}。两法互证，答案可靠。

\aside{注意 $V_{T_2}=\frac{27}{11}\,\mathrm{V}$ 这个数看起来不漂亮，但它是正确的：
本题没有「凑整」的设计。遇到这种答案不必怀疑自己，改用两法互证即可（本题用了变换法与节点法）。
另外，$I=0.818\,\mathrm{mA}$ 若四舍五入成 $0.82\,\mathrm{mA}$，相对误差约 $0.2\%$，作业里可接受，
但中间步骤不要提前取整，否则误差会累积。}

\subhdr{8.5 公式卡：E1 这类题的四步用法}

\callout{\textbf{第 1 步 清场}：找出「与理想电压源并联」和「与理想电流源串联」的元件，删掉
（只删这两类，且源必须是理想的）。\\
\textbf{第 2 步 变换}：见到「电压源 $+$ 串联电阻」，换成 $I_S=U_S/R$ 的电流源\textbf{并联}同一个 $R$；
见到「电流源 $+$ 并联电阻」，换成 $U_S=I_SR$ 的电压源\textbf{串联}同一个 $R$。
箭头尖端对正端。\\
\textbf{第 3 步 合并}：同向并联的电流源相加（反向相减）、并联电阻用 $R_1R_2/(R_1+R_2)$；
串联的电压源相加、串联电阻相加。\\
\textbf{第 4 步 分流（或分压）}：化到「一个源 $+$ 一个分流结构」后，
用 $I_k = I_{\text{总}}\dfrac{R_{\text{并联总}}}{R_k+R_{\text{并联总}}}$ 求待求支路。\\
\textbf{每一步都要重画一遍节点图}——E1 错得最多的一步，是把 $2\,\mathrm{k}\Omega$ 从「串联」误当成「并联」。}

%=============================================================
\hdr{§9}{其余例题、「算不下去」的三处与易错点}

\subhdr{9.1 p59 复习题：用变换求 $I$}

p59 是本节唯一一道要求动笔的复习题：「Find I by using source transformation.」
电路是两支实际电压源并联接一个 $4\,\Omega$ 负载：左支 $40\,\mathrm{V}$ 串 $2\,\Omega$、
中支 $20\,\mathrm{V}$ 串 $2\,\Omega$（两源极性相同，$+$ 都在上），右支 $4\,\Omega$，$I$ 向下。

\begin{center}
\begin{tikzpicture}[scale=0.85,>=Stealth]
  \draw[cwire] (0.5,0) -- (6.5,0);
  \draw[cwire] (0.5,4.4) -- (6.5,4.4);
  % 左支：40V + 2Ω
  \draw[cwire] (0.5,0) -- (0.5,1.2);
  \draw[csrc] (0.5,1.55) circle (0.35);
  \draw[cwire] (0.5,1.9) -- (0.5,2.35);
  \draw[cres] (0.22,2.35) rectangle (0.78,3.35);
  \draw[cwire] (0.5,3.35) -- (0.5,4.4);
  \node[cval,left] at (0.1,1.55) {$40\,\mathrm{V}$};
  \node[cval,right] at (0.85,2.85) {$2\,\Omega$};
  \node[cval,red] at (0.9,2.1) {$+$};
  % 中支：20V + 2Ω
  \draw[cwire] (3.0,0) -- (3.0,1.2);
  \draw[csrc] (3.0,1.55) circle (0.35);
  \draw[cwire] (3.0,1.9) -- (3.0,2.35);
  \draw[cres] (2.72,2.35) rectangle (3.28,3.35);
  \draw[cwire] (3.0,3.35) -- (3.0,4.4);
  \node[cval,left] at (2.6,1.55) {$20\,\mathrm{V}$};
  \node[cval,right] at (3.35,2.85) {$2\,\Omega$};
  % 负载 4Ω 与 I
  \draw[cwire] (6.5,0) -- (6.5,1.4);
  \draw[cres] (6.22,1.4) rectangle (6.78,2.7);
  \draw[cwire] (6.5,2.7) -- (6.5,4.4);
  \node[cval,right] at (6.85,2.05) {$4\,\Omega$};
  \draw[carr,blue] (6.5,4.0) -- (6.5,3.0);
  \node[cval,blue] at (7.05,3.5) {$I$};
  \fill (0.5,4.4) circle (1.4pt); \fill (3.0,4.4) circle (1.4pt); \fill (6.5,4.4) circle (1.4pt);
\end{tikzpicture}
\end{center}

\subhdr{解法一：电源变换（题目要求的方法）}

两支各换成电流源：$40/2 = 20\,\mathrm{A}$ 并 $2\,\Omega$；$20/2 = 10\,\mathrm{A}$ 并 $2\,\Omega$。
两组并联，同向相加：$20+10 = 30\,\mathrm{A}$，$2\,\Omega\parallel2\,\Omega = 1\,\Omega$。
于是化成「$30\,\mathrm{A}$ 并 $1\,\Omega$」接 $4\,\Omega$ 负载，分流：

\[ I = 30\times\frac{1}{1+4} = 6\,\mathrm{A} . \]

\subhdr{解法二：节点法（用来验证）}

取底边为 $0$、上边为 $V$，

\[ \frac{40-V}{2}+\frac{20-V}{2} = \frac{V}{4}
   \;\Longrightarrow\; 40-\tfrac{1}{2}V+20-\tfrac{1}{2}V = \tfrac{1}{4}V
   \;\Longrightarrow\; 120 = 5V
   \;\Longrightarrow\; V = 24\,\mathrm{V}, \]

故 $I = 24/4 = 6\,\mathrm{A}$。两法一致。

\warnbox{\textbf{什么时候用哪一种}：待求量集中在\textbf{一条支路}、而电路本身又要反复化到很简时，
变换法更快（本题 $3$ 步）；节点少、支路多、尤其含受控源时，节点法更稳（本题 $1$ 步）。
\textbf{会两种}的意义在于：考场上一法算出的数可以用另一法在 30 秒内验完。}

\subhdr{9.2 三处「算不下去」的地方，以及为什么算不下去}

\gapitem{（1）E1 为什么不能一上来就用串并联}{p53 的电路有四个节点、两条理想源支路，
「串」与「并」的定义（同一电流／同一电压）在这里都不满足：
$2\,\mathrm{k}\Omega$ 与 $2\,\mathrm{mA}$ 串在一起，它们的电流相同但电压被源钉住，
既不是纯串联也不是纯并联。必须先做「无影响剔除」（§8.1.1）与「等效变换」（§7），
把\textbf{非串并联结构化成串并联结构}，之后才谈得上串并联公式。}

\gapitem{（2）为什么 p52 的「$2\,\mathrm{A}$」不算证明}{见 §7.3。它只验了一个负载点
（$3\,\Omega$）。电源变换是一个「对任意负载都成立」的命题，验证一个点只能说明它「可能」成立。
真正的证明是端口方程的恒等（§7.3 三行）。这不是吹毛求疵：
把 $3\,\Omega$ 换成 $3\,\mathrm{k}\Omega$ 时，两侧是否仍相等？用端口方程一眼可知相等，
用「再验一个点」的办法就要再算一遍。}

\gapitem{（3）为什么变换后不能算内部功率}{见 §7.4 的表：同一对等效模型在空载时，
电压源侧内部损耗 $0\,\mathrm{W}$、电流源侧 $50\,\mathrm{W}$。等效保的是端口，
内部电压、电流、功率全都可以不同。p47 的「For outside not for inside」就是这句话。}

\subhdr{9.3 易错点八条（每条都指回某页某步）}

\par\vspace{2pt}
\noindent\begin{tabular}{@{}p{1.2cm}p{5.4cm}p{7.9cm}@{}}
\tabtop
\thead{序} & \thead{错法} & \thead{指回哪里、怎么改}
\\
\tabmid
（1） & 把参考方向当成真实方向 & p24 第（3）步：算出负数时，改的是「方向」，不是数值。
报答案要写「电流 $2\,\mathrm{A}$，方向与图上箭头相反」
\\
（2） & 非关联方向忘了加负号 & p25 的空图 + p27 三栏表。用 §4.3 的四步法过一遍
\\
（3） & 用 $P=UI$ 而不看方向 & p37。$p=ui$ 只有\textbf{关联}方向下才是吸收功率（§5.1）
\\
（4） & 把 p39 的并联当成 p42 的串联 & 两页元件值一模一样（$3\,\mathrm{A}$、$1\,\Omega$、$3\,\mathrm{V}$），
拓扑却不同。看电阻是竖的还是横的（§5.3）
\\
（5） & 以为「负载增大」是电阻变大 & p10。反了：并联支路增多、等效电阻变小（§1.4）
\\
（6） & 以为与理想电压源并联的电阻会影响电路 & p54 规则（1）。它只改变源自己输出的电流（§6.1.1）
\\
（7） & 变换时把电流源箭头画反 & p48--p50。箭头尖端对正端，否则短路电流反号（§7.2）
\\
（8） & 合并并联电阻时漏掉一支 & p58 的 $\frac{2}{3}\,\mathrm{k}\Omega$。用 $R_1R_2/(R_1+R_2)$ 手算一遍，
               别心算（§8.3）
\tabbot
\end{tabular}

%=============================================================
\hdr{§10}{缺口清单}

下面每一条都按「类型｜页码｜课件英文原文短语｜缺了什么｜补在哪一节」五位一体记录。
共 \textbf{16} 条，其中机制缺口 $6$ 条、代数跳步 $3$ 条、答案形态 $4$ 条、拼接缺口 $1$ 条、
适用区间／前提 $1$ 条、表述（形式改写）缺口 $1$ 条。

\subhdr{10.1 机制缺口（6 条）}

\gapitem{机制缺口}{p5｜「Circuit: A closed path or loop around which current flows.」｜
只说「闭合」，没说为什么必须闭合——若不断言「稳态下电荷不堆积」，就无法解释 p15 的「开路电流恒为零」
以及对「闭合」二字的必然性｜§1.1 用四个部件的能量流与截面电荷守恒补证}

\gapitem{机制缺口}{p2、p4｜「1.1 Function and components」；全册 60 页\textbf{没有出现} model 一词｜
课程与文件名都写「电路模型」，正文却从未定义「模型」，而后面全部公式都建立在「元件只露出两个端钮」这个前提上｜
§1.2 补集总假设（含适用条件与失效情形）}

\gapitem{机制缺口}{p22｜「Notice: Three Representations for Voltage Reference Direction」｜
三种写法只并排放出，没有给\textbf{互译规则}，也没提 $U_{AB}=-U_{BA}$｜
§3.2 补三条互译规则与双下标反对称性；§2.4 给出 $U_{ab}=-U_{ba}$ 的来历}

\gapitem{机制缺口}{p25｜「$U = RI$」（该页矢量图元里既无电流箭头、亦无正负极性）｜
缺了本册最关键的一句：\textbf{非关联参考方向下 $u=-Ri$}。这条不补，p27--p30 三栏表为何都不带负号、
p37 的功率判据为何要谈方向，全都无法解释｜§4.1 给定义、定理与完整证明；
§4.3 给四步判符号法}

\gapitem{机制缺口}{p27｜「$u = -e = L\,\mathrm{d}i/\mathrm{d}t$」｜
一行里压了两件事：物理定律 $u=L\,\mathrm{d}i/\mathrm{d}t$，与 $e$ 的参考极性画反导致的 $u=-e$。
读者容易以为负号是物理的，其实是惯例的｜§4.2.1 拆成两步，并用楞次定律说明为何 $e$ 要反着画}

\gapitem{机制缺口}{p44、p45｜只有 $U$--$I$／$I$--$U$ 两条特性曲线与「$U\ne$ const.」「$I\ne$ const.」｜
没有写出实际源的端口方程 $U=U_S-IR_0$、$I=I_S-U/R_0$，也就无从说明斜率为什么是 $-R_0$｜
§6.3、§6.4 补两式与三截距；这条是 §7 等效变换证明的前置}

\subhdr{10.2 代数跳步（3 条）}

\gapitem{代数跳步}{p10｜「It means the load consumes more electric power, that is, the load current
increases.」｜把「功率增大」与「电流增大」挤成同义句，省略了因果链与前提（端电压不变），
读者会疑心是循环定义｜§1.3 拆成四步推理，并代 $U_S=10\,\mathrm{V}$、$R_1=R_2=2\,\Omega$ 算出 $5\,\mathrm{A}\to10\,\mathrm{A}$}

\gapitem{代数跳步}{p16｜「Short-circuit current is the largest current which may be drawn from the source.」｜
没有给 $I_{\mathrm{SC}}=U_S/R_0$，也没有给闭合状态的 $I=U_S/(R_0+R)$ 与 $U=U_SR/(R_0+R)$｜
§1.6 合成三态表，并代 $10\,\mathrm{V}$／$0.5\,\Omega$／$4.5\,\Omega$ 算出 $2\,\mathrm{A}$、$9\,\mathrm{V}$、短路 $20\,\mathrm{A}$／$200\,\mathrm{W}$}

\gapitem{代数跳步}{p18｜「$V_a=U_{ad}=10\,\mathrm{V}$；$V_b=U_{bd}=6\,\mathrm{V}$；$V_c=U_{cd}=3\,\mathrm{V}$」｜
三个数直接落地，中间省略了回路电流 $1\,\mathrm{A}$ 与逐段压降 $4/3/3\,\mathrm{V}$；
更没交代「换参考点后怎么办」｜§2.2 补四步演算；§2.3 补参考点平移定理并给出 $c$ 作参考点的对照表}

\subhdr{10.3 答案形态（4 条）}

\gapitem{答案形态}{p11｜「（3）Is the circuit a power circuit or a signal circuit?」与
第（4）问的测温电路同样有问无答｜四个提问全部没有答案，而第（4）问恰恰是判据最容易含糊的一处｜
§1.3 逐条作答，并说明第（4）问依据题干而非图面}

\gapitem{答案形态}{p19、p20｜「Find $V_A$ when S is open.」／「when S is closed.」｜
两页只有图与题面，\textbf{一个答案也没有}｜§2.5 全程演算：断开 $-4\,\mathrm{V}$、闭合 $+2\,\mathrm{V}$，
并给出 $V_M=-7.2\,\mathrm{V}$ 的回代校验}

\gapitem{答案形态}{p43｜「Does the current source in the circuit always provide power?」｜
没有答案，而正确答案是「不总是」，需要反接 $3\,\mathrm{V}$ 源才能演示｜
§5.3.1 给出反接后的完整算例：同一个 $3\,\mathrm{A}$ 源由发出 $9\,\mathrm{W}$ 变为吸收 $9\,\mathrm{W}$}

\gapitem{答案形态}{p51、p52（第五问）｜「Why can such a transformation be made?」只用一个 $3\,\Omega$
负载上两侧都是 $2\,\mathrm{A}$ 的演示代替回答｜一个负载点不是证明｜
§7.3 给端口方程恒等的三行证明，并补「两点定一线」的第二个证法；§9.2（2）说明为何演示不算证明}

\subhdr{10.4 拼接缺口、适用区间与表述缺口（3 条）}

\gapitem{拼接缺口}{p41 与 p42｜两页的英文只有同一行标题「Example: Find out which one supply
power, which one absorb power?」，正文对拓扑改动\textbf{一字未提}｜同三个元件值（$3\,\mathrm{A}$、$1\,\Omega$、
$3\,\mathrm{V}$）从「三支并联」变成「一条无分支串联回路」｜按矢量图元核对：p39 的电阻矩形
$20\times64$ pt 为竖向，p42 的 $61\times20$ pt 为横向｜§5.3 指明切换并分别算到底（$1\,\Omega$ 在串联时把
电流源端电压从 $3\,\mathrm{V}$ 抬到 $6\,\mathrm{V}$）}

\gapitem{适用区间／前提}{p54｜「No influence on external circuits」｜
两条「无影响」规则只在源为\textbf{理想}时成立；且第一条会改变理想源自己输出的电流，
「无影响」指的是对外电路而非对源本身｜§6.1.1、§6.2 补两条推论与证明；§8.1.1 落到 E1 的两处剔除上}

\gapitem{表述缺口（形式改写）}{p47｜「It can not only be used to calculate the parameters of the
source itself.」｜与它上一条「For outside not for inside」互相矛盾（not only 把意思写反了）｜
§7.1 按上下文取「只能算外电路」，§7.4 用空载内部损耗 $0\,\mathrm{W}$ 对 $50\,\mathrm{W}$ 的算例
说明「内部不可算」到底是多大差别}

\subhdr{10.5 类型分布与「本册的系统性问题」}

\par\vspace{2pt}
\noindent\begin{tabular}{@{}p{4.2cm}p{1.8cm}p{10.2cm}@{}}
\tabtop
\thead{缺口类型} & \thead{条数} & \thead{集中在哪些页}
\\
\tabmid
机制缺口 & 6 & p2--p5（定义）、p22（三写法）、p25--p27（负号）、p44--p45（实际源方程）
\\
代数跳步 & 3 & p10、p16、p18（都是「结论先行、中间步骤省略」）
\\
答案形态 & 4 & p19--p20、p43、p51--p52（有问无答或答非所问）
\\
拼接缺口 & 1 & p41--p42（拓扑悄悄改变）
\\
适用区间／前提 & 1 & p54（理想与实际的边界）
\\
表述缺口 & 1 & p47（not only 写反）
\tabbot
\end{tabular}

\callout{\textbf{本册的系统性问题（一句话）}：课件的写法是「先给结论、后给演示」，
所以\textbf{机制缺口（6 条）与答案形态（4 条）占了三分之二}。
它把「该讲的那一步」默认读者会补，而本讲义补的正是这些步：
p25 缺一句符号规则、p44 缺一个方程、p51 缺一个证明、p19 缺一个答案。
复习时若只背结论，上面任意一条都会在换一个电路时立刻暴露。}

%=============================================================
\hdr{§11}{三句话总结}

\callout{
\textbf{第一句}：电位要相对参考点才有数，电压是两点之差、与参考点无关；
所以题目问「某点电位」就先定参考点，问「两点电压」就随便取，两者差的就是那一次平移（$V'_x=V_x+U_{ed}$）。\\
\textbf{第二句}：电流进电压的「$+$」端就叫关联，关联时 $u=Ri$、$u=L\,\mathrm{d}i/\mathrm{d}t$、
$i=C\,\mathrm{d}u/\mathrm{d}t$、$p=ui$ 且 $p>0$ 为吸收；电流出「$+$」端就全部加负号。
负号不是新物理，是记号与物理方向不一致的记账费。\\
\textbf{第三句}：等效变换只换「端口方程」，$R'=R$、$I_S=U_S/R$ 的根据就是
$U=U_S-IR$ 与 $I=I_S-U/R'$ 是同一个方程；它保外不保内，
所以能用来算外电路，不能用来算电源内部的电流与损耗。}

%=============================================================
\hdr{§12}{与课件的页码对照表}

\par\vspace{2pt}
\noindent\begin{tabular}{@{}p{1.8cm}p{5.4cm}p{2.4cm}p{6.2cm}@{}}
\tabtop
\thead{课件页} & \thead{页上内容} & \thead{本讲义节} & \thead{处理方式}
\\
\tabmid
p1 & Contents（七章目录） & §0.1 & 只取「本册在全书中的位置」
\\
p2--p4 & 1.1 纲目、四项内容 & §0.1、§1 & 原样转述
\\
p5 & 电路定义 + 实物／符号图 & §1.1 & 原样 + 补「为什么必须闭合」
\\
p6 & 按功能分类的框架图＋提问 & §1.3 & 原样 + 补判据
\\
p7--p8 & 电力电路／信号电路的组成与关注点 & §1.3 & 原样（表内四行全引）
\\
p9 & 三个组成部分 + 手电筒实例 & §1.4 & 原样
\\
p10 & 「Load Increase」的含义 & §1.3 & 补因果链四步 + 代 $10\,\mathrm{V}$／$2\,\Omega$ 的数
\\
p11 & 四个提问 & §1.3 & \textbf{补全部答案}（§10）
\\
p12 & 1.2 基本量清单（六行名词） & §1.5 & 补「单位」与「本册首次用到处」两列
\\
p13--p14 & 三种状态、闭合电路 & §1.6 & 补 $I$、$U$ 两式（§10）
\\
p15 & 开路：$I=0$、$U=U_{\mathrm{OC}}=U_S$、$P=0$ & §1.6 & 原样 + 补「为什么」
\\
p16 & 短路：$U=0$、$I$ 任意 & §1.6 & \textbf{补 $I_{\mathrm{SC}}=U_S/R_0$} 与代数值
\tabbot
\end{tabular}

\par\vspace{10pt}
\noindent\begin{tabular}{@{}p{1.8cm}p{5.4cm}p{2.4cm}p{6.2cm}@{}}
\tabtop
\thead{课件页} & \thead{页上内容} & \thead{本讲义节} & \thead{处理方式}
\\
\tabmid
p17--p18 & 参考点例：三问与三个答案 & §2.1、§2.2 & \textbf{补中间步骤}（§10）
\\
p19--p20 & $V_A$ 两问（S 断开／闭合） & §2.5 & \textbf{补全部答案}（§10）
\\
p21 & 关注的方向与「参考方向＝假定」 & §3.1 & 原样 + 点明 assumed 的分量
\\
p22 & 电压参考方向三写法 & §3.2 & 补互译规则与 $U_{AB}=-U_{BA}$（§10）
\\
p23--p24 & Example 1（已知真实方向）／Example 2（先假定） & §3.1、§3.3 & 原样 + 补「为什么假定不会错」
\\
p25 & Example 3：$U=RI$（无方向标注） & §4.1 & \textbf{补关联定义、定理与证明}（§10）
\\
p26 & 理想元件分类树＋提问 & §4.4 & 提问并入 §5 的判据
\\
p27--p30 & 三栏表（R／L／C）与四页结论句 & §4.2 & 逐栏判读参考方向 + 拆 $u=-e$ 的两个负号
\\
p31--p32 & 忆阻器（$1971$ Chua、$2008$ HP）与四象限图 & §4.2.2 & 四象限关系只作背景，（4）（5）两栏课件留空
\tabbot
\end{tabular}

\par\vspace{10pt}
\noindent\begin{tabular}{@{}p{1.8cm}p{5.4cm}p{2.4cm}p{6.2cm}@{}}
\tabtop
\thead{课件页} & \thead{页上内容} & \thead{本讲义节} & \thead{处理方式}
\\
\tabmid
p33--p34 & 理想电压源特性 + $10\,\mathrm{V}$ 例 & §6.1 & 原样 + 补两条禁令
\\
p35--p36 & 理想电流源特性 + $1\,\mathrm{A}$ 例 & §6.2 & 原样 + 补两条禁令与对偶
\\
p37 & $P=UI$ 与吸收／发出判据 & §5.1 & \textbf{补三处漏洞}（§10）
\\
p38--p41 & 3A／3V 并联组（$R$ 无、$1$、$0.5$、$2\,\Omega$） & §5.2 & \textbf{补四行完整计算表}（课件无结论）
\\
p42--p43 & 改为串联；问「电流源是否总发出」 & §5.3 & 指出拓扑切换 + \textbf{补答案}（§10）
\\
p44--p45 & 理想／实际电压源、电流源对照 & §6.3、§6.4 & \textbf{补两个端口方程}
\tabbot
\end{tabular}

\par\vspace{10pt}
\noindent\begin{tabular}{@{}p{1.8cm}p{5.4cm}p{2.4cm}p{6.2cm}@{}}
\tabtop
\thead{课件页} & \thead{页上内容} & \thead{本讲义节} & \thead{处理方式}
\\
\tabmid
p46 & 五问中的（1）（2） & §7.1 & 原样
\\
p47 & 五问中的（3）：两条条件 & §7.1、§7.4 & 指出 not only 的表述矛盾 + 补定量版
\\
p48 & $R'=R$、$I_S=U_S/R$ 与两个模型 & §7.2 & 补「箭头尖端对正端」的方向规则
\\
p49--p50 & $10\,\mathrm{V}$ 串 $2\,\Omega\to5\,\mathrm{A}$ 并 $2\,\Omega$ & §7.2 & 原样复核数字
\\
p51--p52 & 接 $3\,\Omega$ 负载两侧均 $2\,\mathrm{A}$ & §7.3 & 原样 + \textbf{补证明}（§10）
\\
p53--p58 & 例 E1 五步化简 & §8 & 复原节点表 + 补两处剔除的证明 + \textbf{补最终答案 $\frac{9}{11}\,\mathrm{mA}$}
\\
p59 & 复习题：求 $I$ & §9.1 & \textbf{补答案 $6\,\mathrm{A}$}（变换法与节点法各一遍）
\\
p60 & Review 提纲（三大部分、五个小问） & §11 & 按同一框架写三句话总结
\tabbot
\end{tabular}

\flowerrule

{\small\color{inkgray}
本讲义依据课件 \texttt{Chapter one -1-2026中白上传.pdf}（共 60 页）编写；
连接关系按 PDF 矢量图元坐标逐页复原，数字全部经独立计算复核
（$V_a,V_b,V_c$；$I_{\mathrm{SC}}$；四张 3A/3V 表的 $\sum p=0$；$V_A=-4\,\mathrm{V}$ 与 $+2\,\mathrm{V}$；
E1 的 $\frac{9}{11}\,\mathrm{mA}$ 与节点法互证；p59 的 $6\,\mathrm{A}$ 两法互证）。
标注 [教材补充] 者为本讲义为补足逻辑而加、课件页面上没有的内容。}

\end{document}







